\section{AIME I 2024}
        \begin{enumerate}
        	\item (\textit{AIME I 2024}) Every morning Aya goes for a $9$-kilometer-long walk and stops at a coffee shop afterwards. When she walks at a constant speed of $s$ kilometers per hour, the walk takes her $4$ hours, including $t$ minutes spent in the coffee shop. When she walks at $s + 2$ kilometers per hour, the walk takes her $2$ hours and $24$ minutes, including $t$ minutes spent in the coffee shop. Suppose Aya walks at $s + \frac{1}{2}$
kilometers per hour. Find the number of minutes the walk takes her, including the $t$ minutes spent in the coffee shop.


 
 \textbf{Answer: }204



\textbf{Solution 1}

$\frac{9}{s} + t = 4$ in hours and $\frac{9}{s+2} + t = 2.4$ in hours.

 Subtracting the second equation from the first, we get, 

 $\frac{9}{s} - \frac{9}{s+2} = 1.6$

 Multiplying by $(s)(s+2)$, we get 

 $9s+18-9s=18=1.6s^{2} + 3.2s$

 Multiplying by 5/2 on both sides, we get

 $0 = 4s^{2} + 8s - 45$

 Factoring gives us 

 $(2s-5)(2s+9) = 0$, of which the solution we want is $s=2.5$.

 Substituting this back to the first equation, we can find that $t = 0.4$ hours.

 Lastly, $s + \frac{1}{2} = 3$ kilometers per hour, so

 $\frac{9}{3} + 0.4 = 3.4$ hours, or $\framebox{204}$ minutes

 -Failure.net

 





\textbf{Solution 2}

The amount of hours spent while walking on the first travel is $\frac{240-t}{60}$. Thus, we have the equation $(240-t)(s) = 540$, and by the same logic, the second equation yields $(144-t)(s+2) = 540$. We have $240s-st = 540$, and $288+144s-2t-st = 540$. We subtract the two equations to get $96s+2t-288 = 0$, so we have $48s+t = 144$, so $t = 144-48s$, and now we have $(96+48s)(s) = 540$. The numerator of $s$ must evenly divide 540, however, $s$ must be less than 3. We can guess that $s = 2.5$. Now, $2.5+0.5 = 3$. Taking $\frac{9}{3} = 3$, we find that it will take three hours for the 9 kilometers to be traveled. The t minutes spent at the coffee shop can be written as $144-48(2.5)$, so t = 24. $180 + 24 =  \framebox{204}$. -sepehr2010

 


 




	\item (\textit{AIME I 2024}) There exist real numbers $x$ and $y$, both greater than 1, such that $\log_x(y^x) = \log_y(x^{4y}) = 10.$ Find $xy$.


 
 \textbf{Answer: }025



\textbf{Solution 1}

By properties of logarithms, we can simplify the given equation to $x\log_xy=4y\log_yx=10$. Let us break this into two separate equations: 

 
 \[x\log_xy=10\]
 \[4y\log_yx=10.\]We multiply the two equations to get: 
 \[4xy\left(\log_xy\log_yx\right)=100.\]

 Also by properties of logarithms, we know that $\log_ab\cdot\log_ba=1$; thus, $\log_xy\cdot\log_yx=1$. Therefore, our equation simplifies to: 

 
 \[4xy=100\implies xy=\boxed{025}.\]

 





\textbf{Solution 2}

Convert the two equations into exponents:

 
 \[x^{10}=y^x~(1)\]
 \[y^{10}=x^{4y}~(2).\]

 Take $(1)$ to the power of $\frac{1}{x}$:

 
 \[x^{\frac{10}{x}}=y.\]

 Plug this into $(2)$:

 
 \[x^{(\frac{10}{x})(10)}=x^{4(x^{\frac{10}{x}})}\]
 \[{\frac{100}{x}}={4x^{\frac{10}{x}}}\]
 \[{\frac{25}{x}}={x^{\frac{10}{x}}}=y,\]

 So $xy=\boxed{025}$

 





\textbf{Solution 3}

Similar to solution 2, we have:

 $x^{10}=y^x$ and $y^{10}=x^{4y}$

 Take the tenth root of the first equation to get 

 $x=y^{\frac{x}{10}}$ 

 Substitute into the second equation to get 

 $y^{10}=y^{\frac{4xy}{10}}$

 This means that $10=\frac{4xy}{10}$, or $100=4xy$, meaning that $xy=\boxed{025}$.

 





\textbf{Solution 4}

The same with other solutions, we have obtained $x^{10}=y^x$ and $y^{10}=x^{4y}$. Then, $x^{10}y^{10}=y^xx^{4y}$. So, an obvious solution is to have $x^{10}=x^{4y}$ and $y^{10}=y^{x}$. Solving, we get $x=10$ and $y=2.5$.So $xy = \boxed{025}$.

 Note: This is not a correct solution. Plugging in $x=10$ and $y=2.5$ does not satisfy the equations.

 





\textbf{Solution 5(fakesolve)}

Using the first expression, we see that $x^{10} = y^x$. Now, taking the log of both sides, we get $\log_y(x^{10}) = \log_y(y^x)$. This simplifies to $10 \log_y(x) = x$. This is still equal to the second equation in the problem statement, so $10 \log_y(x) = x = 4y \log_y(x)$. Dividing by $\log_y(x)$ on both sides, we get $x = 4y = 10$. Therefore, $x = 10$ and $y = 2.5$, so $xy = \boxed{025}$.

 





\textbf{Solution 6}

Let 
 \[y=x^a\].We see: 
 \[ax=10\] and 
 \[4x^a/a=10\]which gives rise to 
 \[xy = \boxed{025}\].

 





\textbf{Solution 7 (log rules)}

Split into 2 equations to get 
 \[\log_x\left(y^x\right)=10\] 
 \[\log_y\left(x^{4y}\right)=10\]Now use the logarithm power rule and the change of base rule to obtain 
 \[x*\frac{\log y}{\log x}=10\] 
 \[4y*\frac{\log x}{\log y}=10\] Multiplying both equations to cancel $\log x$ and $\log y$, we get 
 \[4xy = 100\] 
 \[xy = \boxed{025}\]

 


 





\textbf{Solution 8 (Exponential Form)}

We can convert the equality chain into two equations in exponential form, which gives 
 \[x^{10} = y^{x}\] and 
 \[y^{10} = x^{4y}\]To obtain an equation which involves exponents of only one variable, we can raise to first equation to power of 10, and raise the second equation to the power of x, which gives 
 \[x^{100} = y^{10x}\] and 
 \[x^{4xy} = y^{10x}\] Now we can deduce that $x^{4xy} = x^{100}$, which must mean that $4xy = 100$ and the desired value of 
 \[xy = \boxed{025}\]

 


 




	\item (\textit{AIME I 2024}) Alice and Bob play the following game. A stack of $n$ tokens lies before them. The players take turns with Alice going first. On each turn, the player removes $1$ token or $4$ tokens from the stack. The player who removes the last token wins. Find the number of positive integers $n$ less than or equal to $2024$ such that there is a strategy that guarantees that Bob wins, regardless of Alice’s moves.


 
 \textbf{Answer: }809



\textbf{Solution 1}

Let's first try some experimentation. Alice obviously wins if there is one coin. She will just take it and win. If there are 2 remaining, then Alice will take one and then Bob will take one, so Bob wins. If there are $3$, Alice will take $1$, Bob will take one, and Alice will take the final one. If there are $4$, Alice will just remove all $4$ at once. If there are $5$, no matter what Alice does, Bob can take the final coins in one try. Notice that Alice wins if there are $1$, $3$, or $4$ coins left. Bob wins if there are $2$ or $5$ coins left.

 After some thought, you may realize that there is a strategy for Bob. If there is n is a multiple of $5$, then Bob will win. The reason for this is the following: Let's say there are a multiple of $5$ coins remaining in the stack. If Alice takes $1$, Bob will take $4$, and there will still be a multiple of $5$. If Alice takes $4$, Bob will take $1$, and there will still be a multiple of $5$. This process will continue until you get $0$ coins left. For example, let's say there are $205$ coins. No matter what Alice does, Bob can simply just do the complement. After each of them make a turn, there will always be a multiple of $5$ left. This will continue until there are $5$ coins left, and Bob will end up winning.

 After some more experimentation, you'll realize that any number that is congruent to $2$ mod $5$ will also work. This is because Bob can do the same strategy, and when there are $2$ coins left, Alice is forced to take $1$ and Bob takes the final coin. For example, let's say there are $72$ coins. If Alice takes $1$, Bob will take $4$. If Alice takes $4$, Bob will take $1$. So after they each make a turn, the number will always be equal to $2$ mod $5$. Eventually, there will be only $2$ coins remaining, and we've established that Alice will simply take $1$ and Bob will take the final coin.

 So we have to find the number of numbers less than or equal to $2024$ that are either congruent to $0$ mod $5$ or $2$ mod $5$. There are $404$ numbers in the first category: $5, 10, 15, \dots, 2020$. For the second category, there are $405$ numbers. $2, 7, 12, 17, \dots, 2022$. So the answer is $404 + 405 = \boxed{809}$

 





\textbf{Solution 2}

Let's use winning and losing positions, where $W$ marks a win for Alice.

 $1$ coin: $W$

 $2$ coins: $L$

 $3$ coins: $W$

 $4$ coins: $W$

 $5$ coins: $L$

 $6$ coin: $W$

 $7$ coins: $L$

 $8$ coins: $W$

 $9$ coins: $W$

 $10$ coins: $L$

 $11$ coin: $W$

 $12$ coins: $L$

 $13$ coins: $W$

 $14$ coins: $W$

 $15$ coins: $L$

 We can see that losing positions occur when $n$ is congruent to $0, 2 \mod{5}$ and winning positions occur otherwise. As $n$ ranges from $1$ to $2020$, $\frac{2}{5}$ of these values are losing positions where Bob will win. As $n$ ranges from $2021$ to $2024$, $2022$ is the only value where Bob will win. Thus, the answer is $2020\times\frac{2}{5}+1=\boxed{809}$

 





\textbf{Solution 3}

Denote by $A_i$ and $B_i$ Alice's or Bob's $i$th moves, respectively.

 Case 1: $n \equiv 0 \pmod{5}$.

 Bob can always take the strategy that $B_i = 5 - A_i$.This guarantees him to win.

 In this case, the number of $n$ is $\left\lfloor \frac{2024}{5} \right\rfloor = 404$.

 Case 2: $n \equiv 1 \pmod{5}$.

 In this case, consider Alice's following strategy: $A_1 = 1$ and $A_i = 5 - B_{i-1}$ for $i \geq 2$.Thus, under Alice's this strategy, Bob has no way to win.

 Case 3: $n \equiv 4 \pmod{5}$.

 In this case, consider Alice's following strategy: $A_1 = 4$ and $A_i = 5 - B_{i-1}$ for $i \geq 2$.Thus, under Alice's this strategy, Bob has no way to win.

 Case 4: $n \equiv 2 \pmod{5}$.

 Bob can always take the strategy that $B_i = 5 - A_i$.Therefore, after the $\left\lfloor \frac{n}{5} \right\rfloor$th turn, there are two tokens leftover.Therefore, Alice must take 1 in the next turn that leaves the last token on the table.Therefore, Bob can take the last token to win the game.This guarantees him to win.

 In this case, the number of $n$ is $\left\lfloor \frac{2024 - 2}{5} \right\rfloor +1 = 405$.

 Case 5: $n \equiv 3 \pmod{5}$.

 Consider Alice's following strategy: $A_1 = 1$ and $A_i = 5 - B_{i-1}$ for $i \geq 2$.By doing so, there will finally be 2 tokens on the table and Bob moves first. Because Bob has the only choice of taking 1 token, Alice can take the last token and win the game.

 Therefore, in this case, under Alice's this strategy, Bob has no way to win.

 Putting all cases together, the answer is $404 + 405 = \boxed{\textbf{(809) }}$.

 


 





\textbf{Solution 4 (Grundy Values)}

Since the game Alice and Bob play is impartial (the only difference between player 1 and player 2 is that player 1 goes first (note that games like chess are not impartial because each player can only move their own pieces)), we can use the Sprague-Grundy Theorem to solve this problem. We will use induction to calculate the Grundy Values for this game.

 We claim that heaps of size congruent to $0,2 \bmod{5}$ will be in outcome class $\mathcal{P}$ (win for player 2 = Bob), and heaps of size equivalent to $1,3,4 \bmod{5}$ will be in outcome class $\mathcal{N}$ (win for player 1 = Alice). Note that the mex (minimal excludant) of a set of nonnegative integers is the least nonnegative integer not in the set. e.g. mex$(1, 2, 3) = 0$ and mex$(0, 1, 2, 4) = 3$.

 
$\text{heap}(0) = \{\} = *\text{mex}(\emptyset) = 0$

 
$\text{heap}(1) = \{0\} = *\text{mex}(0) = *$

 
$\text{heap}(2) = \{*\} = *\text{mex}(1) = 0$

 
$\text{heap}(3) = \{0\} = *\text{mex}(0) = *$

 
$\text{heap}(4) = \{0, *\} = *\text{mex}(0, 1) = *2$

 
$\text{heap}(5) = \{*, *2\} = *\text{mex}(1, 2) = 0$

 
$\text{heap}(6) = \{0, 0\} = *\text{mex}(0, 0) = *$

 
$\text{heap}(7) = \{*, *\} = *\text{mex}(1, 1) = 0$

 
$\text{heap}(8) = \{*2, 0\} = *\text{mex}(0, 2) = *$

 
$\text{heap}(9) = \{0, *\} = *\text{mex}(0, 1) = *2$

 
$\text{heap}(10) = \{*, *2\} = *\text{mex}(1, 2) = 0$

 
We have proven the base case. We will now prove the inductive hypothesis: If $n \equiv 0 \bmod{5}$, $\text{heap}(n) = 0$, $\text{heap}(n+1) = *$, $\text{heap}(n+2) = 0$, $\text{heap}(n+3) = *$, and $\text{heap}(n+4) = *2$, then $\text{heap}(n+5) = 0$, $\text{heap}(n+6) = *$, $\text{heap}(n+7) = 0$, $\text{heap}(n+8) = *$, and $\text{heap}(n+9) = *2$.

 
$\text{heap}(n+5) = \{\text{heap}(n+1), \text{heap}(n+4)\} = \{*, *2\} = *\text{mex}(1, 2) = 0$

 
$\text{heap}(n+6) = \{\text{heap}(n+2), \text{heap}(n+5)\} = \{0, 0\} = *\text{mex}(0, 0) = *$

 
$\text{heap}(n+7) = \{\text{heap}(n+3), \text{heap}(n+6)\} = \{*, *\} = *\text{mex}(1, 1) = 0$

 
$\text{heap}(n+8) = \{\text{heap}(n+4), \text{heap}(n+7)\} = \{*2, 0\} = *\text{mex}(2, 1) = *$

 
$\text{heap}(n+9) = \{\text{heap}(n+5), \text{heap}(n+8)\} = \{0, *\} = *\text{mex}(0, 1) = *2$

 
We have proven the inductive hypothesis. QED.

 There are $2020*\frac{2}{5}=808$ positive integers congruent to $0,2 \bmod{5}$ between 1 and 2020, and 1 such integer between 2021 and 2024. $808 + 1 = \boxed{809}$.

 





\textbf{Solution 5 (no modular arithmetic)}

We start with $n$ as some of the smaller values. After seeing the first 4 where Bob wins automatically, with trial and error we see that $2, 5, 7,$ and $10$ are spaced alternating in between 2 and 3 apart. This can also be proven with modular arithmetic, but  this is an easier solution for some people. We split them into 2 different sets with common difference 5: {2,7,12 ...} and {5,10,15...}. Counting up all the numbers in each set can be done as follows:

 Set 1 ${2,7,12...}$

 $2024-2=2022$ (because the first term is two)

 $\lfloor \frac{2024}{5} \rfloor = 404$

 Set 2 ${5,10,15}$

 $\lfloor \frac{2024}{5} \rfloor = 404$

 
And because we forgot 2022 we add 1 more.

 $404+404+1=809$

 -Multpi12(Edits would be appreciated)LaTexed by BossLu99

 


 





\textbf{Solution 6 (Easiest, Very Intuitive, Fast)}

The trivial case is all 0 (mod 5)s because for every 1 or 4 Alice chooses Bob can choose the opposite and they always sum to 5. Now we test 1, 2, 3, 4 (mod 5). 4 (mod 5) does not work because Alice would take all 4, 3 (mod 5) does not work because it would go A-B-A, in which Alice wins, 2 (mod 5) works because since Alice has to go first, she has to take 1 then Bob wins, 1 (mod 5) does not work because Alice would just take the remainder of 1. 

 Thus, only 0 (mod 5) and 2 (mod 5)s work. This is 809. 

 




	\item (\textit{AIME I 2024}) Jen enters a lottery by selecting $4$ distinct elements of $S=\{1,2,3,\cdots,9,10\}.$ Then four elements of $S$ are drawn at random. Jen wins a prize if at least two of her numbers were drawn, and wins the grand prize if all four of her numbers are drawn. The probability that Jen wins the grand prize given that Jen wins a prize is $\frac{m}{n}$ where $m$ and $n$ are relatively prime positive integers. Find $m+n$.


 
 \textbf{Answer: }116



\textbf{Solution 1}

This is a conditional probability problem. Bayes' Theorem states that 
 \[P(A|B)=\dfrac{P(B|A)\cdot P(A)}{P(B)}\]

 -- in other words, the probability of $A$ given $B$ is equal to the probability of $B$ given $A$ times the probability of $A$ divided by the probability of $B$. In our case, $A$ represents the probability of winning the grand prize, and $B$ represents the probability of winning a prize. Clearly, $P(B|A)=1$, since by winning the grand prize you automatically win a prize. Thus, we want to find $\dfrac{P(A)}{P(B)}$. 

 Let us calculate the probability of winning a prize. We do this through casework: how many of Jen's drawn numbers match the lottery's drawn numbers? 

 To win a prize, Jen must draw at least $2$ numbers identical to the lottery. Thus, our cases are drawing $2$, $3$, or $4$ numbers identical. 

 Let us first calculate the number of ways to draw exactly $2$ identical numbers to the lottery. Let Jen choose the numbers $a$, $b$, $c$, and $d$; we have $\dbinom42$ ways to choose which $2$ of these $4$ numbers are identical to the lottery. We have now determined $2$ of the $4$ numbers drawn in the lottery; since the other $2$ numbers Jen chose can not be chosen by the lottery, the lottery now has $10-2-2=6$ numbers to choose the last $2$ numbers from. Thus, this case is $\dbinom62$, so this case yields $\dbinom42\dbinom62=6\cdot15=90$ possibilities. 

 Next, let us calculate the number of ways to draw exactly $3$ identical numbers to the lottery. Again, let Jen choose $a$, $b$, $c$, and $d$. This time, we have $\dbinom43$ ways to choose the identical numbers and again $6$ numbers left for the lottery to choose from; however, since $3$ of the lottery's numbers have already been determined, the lottery only needs to choose $1$ more number, so this is $\dbinom61$. This case yields $\dbinom43\dbinom61=4\cdot6=24$. 

 Finally, let us calculate the number of ways to all $4$ numbers matching. There is actually just one way for this to happen, as $\dbinom44=1$. 

 In total, we have $90+24+1=115$ ways to win a prize. The lottery has $\dbinom{10}4=210$ possible combinations to draw, so the probability of winning a prize is $\dfrac{115}{210}$. There is actually no need to simplify it or even evaluate $\dbinom{10}4$ or actually even know that it has to be $\dbinom{10}4$; it suffices to call it $a$ or some other variable, as it will cancel out later. However, let us just go through with this. The probability of winning a prize is $\dfrac{115}{210}$. Note that the probability of winning a grand prize is just matching all $4$ numbers, which we already calculated to have $1$ possibility and thus have probability $\dfrac1{210}$. Thus, our answer is $\dfrac{\frac1{210}}{\frac{115}{210}}=\dfrac1{115}$. Therefore, our answer is $1+115=\boxed{116}$. 

 





\textbf{Solution 2(quick)}

One may also use complimentary counting as a shortcut to calculate the probability of winning a prize, in which the cases are that either one number is shared or no numbers are shared. There are $4 \cdot { {10 - 4} \choose {4-1}} = 4 \cdot 20 = 80$ ways to choose the former and ${{10-4} \choose 4} = 15$ ways for the latter. Therefore, there are $95$ ways to NOT choose a prize, so there are $210-95 = 115$ ways to choose a prize, and the answer follows.

 - \href{https://artofproblemsolving.com/wiki/index.php/User:Spectraldragon8}{spectraldragon8}

 





\textbf{Solution 3}

For getting all $4$ right, there is only $1$ way.

 For getting $3$ right, there is $\dbinom43$ multiplied by $\dbinom61$ = $24$ ways.

 For getting $2$ right, there is $\dbinom42$ multiplied by $\dbinom62$ = $90$ ways.

 $\frac{1}{1+24+90}$  = $\frac{1}{115}$ 

 Therefore, the answer is $1+115 = \boxed{116}$

 


 




	\item (\textit{AIME I 2024}) Rectangle $ABCD$ has dimensions $AB = 107$ and $BC = 16$, and rectangle $EFGH$ has dimensions $EF = 184$ and $FG = 17$. Points $D$, $E$, $C$, and $F$ lie on line $DF$ in that order, and $A$ and $H$ lie on opposite sides of line $DF$, as shown. Points $A$, $D$, $H$, and $G$ lie on a common circle. Find $CE$.


 
                    \begin{figure}[H]
                        \centering
                        \includegraphics[width=0.5\textwidth]{Problem/AIME_2024_I_5_Problem_Figure_1.png}
                    \end{figure}
                    


 
 \textbf{Answer: }104



\textbf{Solution 1}

We use simple geometry to solve this problem. 

 
                    \begin{figure}[H]
                        \centering
                        \includegraphics[width=0.5\textwidth]{Solution/AIME_2024_I_5_Solution_Figure_1.png}
                    \end{figure}
                    

 We are given that $A$, $D$, $H$, and $G$ are concyclic; call the circle that they all pass through circle $\omega$ with center $O$. We know that, given any chord on a circle, the perpendicular bisector to the chord passes through the center; thus, given two chords, taking the intersection of their perpendicular bisectors gives the center. We therefore consider chords $HG$ and $AD$ and take the midpoints of $HG$ and $AD$ to be $P$ and $Q$, respectively. 

 
                    \begin{figure}[H]
                        \centering
                        \includegraphics[width=0.5\textwidth]{Solution/AIME_2024_I_5_Solution_Figure_2.png}
                    \end{figure}
                    

 We could draw the circumcircle, but actually it does not matter for our solution; all that matters is that $OA=OH=r$, where $r$ is the circumradius. 

 By the Pythagorean Theorem, $OQ^2+QA^2=OA^2$. Also, $OP^2+PH^2=OH^2$. We know that $OQ=DE+HP$, and $HP=\dfrac{184}2=92$; $QA=\dfrac{16}2=8$; $OP=DQ+HE=8+17=25$; and finally, $PH=92$. Let $DE=x$. We now know that $OA^2=(x+92)^2+8^2$ and $OH^2=25^2+92^2$. Recall that $OA=OH$; thus, $OA^2=OH^2$. We solve for $x$: 

 \begin{align*}(x+92)^2+8^2&=25^2+92^2 \\(x+92)^2&=625+(100-8)^2-8^2 \\&=625+10000-1600+64-64 \\&=9025 \\x+92&=95 \\x&=3. \\\end{align*}


 The question asks for $CE$, which is $CD-x=107-3=\boxed{104}$.

 


 





\textbf{Solution 2}

Suppose $DE=x$. Extend $AD$ and $GH$ until they meet at $P$. From the \href{https://artofproblemsolving.com/wiki/index.php?title=Power\_of\_a\_Point\_Theorem}{Power of a Point Theorem}, we have $(PH)(PG)=(PD)(PA)$. Substituting in these values, we get $(x)(x+184)=(17)(33)=561$. We can use guess and check to find that $x=3$, so $EC=\boxed{104}$.
                    \begin{figure}[H]
                        \centering
                        \includegraphics[width=0.5\textwidth]{Solution/AIME_2024_I_5_Solution_Figure_3.png}
                    \end{figure}
                    

 rabbit47 - quadratic actually factors as (x+187)(x-3)=0, from which x=3

 





\textbf{Solution 3}

We find that 
 \[\angle GAB = 90-\angle DAG = 90 - (180 - \angle GHD) = \angle DHE.\]

 Let $x = DE$ and $T = FG \cap AB$. By similar triangles $\triangle DHE \sim \triangle GAT$ we have $\frac{DE}{EH} = \frac{GT}{AT}$. Substituting lengths we have $\frac{x}{17} = \frac{16 + 17}{184 + x}.$ Solving, we find $x = 3$ and thus $CE = 107 - 3 = \boxed{104}.$~AtharvNaphade ~coolruler ~eevee9406

 





\textbf{Solution 4}

One liner: $107-\sqrt{92^2+25^2-8^2}+92=\boxed{104}$

 \textbf{Explanation} Let $OP$ intersect $DF$ at $T$ (using the same diagram as Solution 2).

 The formula calculates the distance from $O$ to $H$ (or $G$), $\sqrt{92^2+25^2}$, then shifts it to $OD$ and the finds the distance from $O$ to $Q$, $\sqrt{92^2+25^2-8^2}$. $107$ minus that gives $CT$, and when added to $92$, half of $FE=TE$, gives $CT+TE=CE$

 





\textbf{Solution 5}

Let $\angle{DHE} = \theta.$ This means that $DE = 17\tan{\theta}.$ Since quadrilateral $ADHG$ is cyclic, $\angle{DAG} = 180 - \angle{DHG} = 90 - \theta.$

 Let $X = AG \cap DF.$ Then, $\Delta DXA \sim \Delta FXG,$ with side ratio $16:17.$ Also, since $\angle{DAG} = 90 - \theta, \angle{DXA} = \angle{FXG} = \theta.$ Using the similar triangles, we have $\tan{\theta} = \frac{16}{DX} = \frac{17}{FX}$ and $DX + FX = DE + EF = 17\tan{\theta} + 184.$

 Since we want $CE = CD - DE = 107 - 17\tan{\theta},$ we only need to solve for $\tan{\theta}$ in this system of equations. Solving yields $\tan{\theta} = \frac{3}{17},$ so $CE = \boxed{104.}$

 





\textbf{Solution 6}

Using a ruler (also acting as a straight edge), draw the figure to scale with one unit = 1mm. With a compass, draw circles until you get one such that $A,D,H,G$ are on the edge of the drawn circle. From here, measuring with your ruler should give $CE = \boxed{104.}$

 Note: 1 mm is probably the best unit to use here just for convenience (drawing all required parts of the figure fits into a normal-sized scrap paper 8.5 x 11); also all lines can be drawn with a standard 12-inch ruler 

 





\textbf{Solution 7(Ptolemy's Theorem)}

Since ADHG is cyclic, AH*DG=AD*HG+DH*AG,let DE be x, we represent the equation in terms of x,and you get $x=3, 107-3= \boxed{104.}$

 Ps: extremely time consuming, do not use on the test

 \textbf{Fast Video Solution by Do Math or Go Home} 




	\item (\textit{AIME I 2024}) Consider the paths of length $16$ that follow the lines from the lower left corner to the upper right corner on an $8\times 8$ grid. Find the number of such paths that change direction exactly four times, as in the examples shown below.


 
                    \begin{figure}[H]
                        \centering
                        \includegraphics[width=0.5\textwidth]{Problem/AIME_2024_I_6_Problem_Figure_1.png}
                    \end{figure}
                    


 
 \textbf{Answer: }294



\textbf{Solution 1}

We divide the path into eight “$R$” movements and eight “$U$” movements. Five sections of alternating $RURUR$ or $URURU$ are necessary in order to make four “turns" because each transition from $R$ to $U$ or from $U$ to $R$ is one turn. Without loss of generality, we use the first case and multiply by $2$ (because we can convert each path in the first case to one in the second case by turning $R$'s to $U$'s and $U$'s to $R$'s).

 
For $U$, let $a$ be the length of the first block of $U$'s and $b$ be the second block of $U$'s. There are eight $U$'s in total, so we have to have $a+b=8$. We have seven ordered pairs of such positive integers $(a,b)$, specifically $(1,7),(2,6),\ldots,(6,2),(7,1)$.

 For $R$, we subtract $1$ from each block of $R$'s (because each block of $R$'s must have length at least $1$, and to use stars and bars we have to have minimum length $0$). From here, we use stars and bars to get ${7 \choose 5}=21$.

 
Thus our answer is $7\cdot21\cdot2=\boxed{294}$.

 





\textbf{Solution 2}

Draw a few examples of the path. However, notice one thing in common - if the path starts going up, there will be 3 "segments" where the path goes up, and two horizontal "segments." Similarly, if the path starts going horizontally, we will have three horizontal segments and two vertical segments. Those two cases are symmetric, so we only need to consider one. If our path starts going up, by stars and bars, we can have $\binom{7}{2}$ ways to split the 8 up's into 3 lengths, and have $\binom{7}{1}$ to split the 8-horizontals into 2 lengths. We multiply them together, and multiply by 2 for symmetry, giving us $2*\binom{7}{2}*\binom{7}{1}=294.$

 





\textbf{Solution 3}

Notice that the $RURUR$ case and the $URURU$ case is symmetrical. WLOG, let's consider the RURUR case.

 Now notice that there is a one-to-one correspondence between this problem and the number of ways to distribute 8 balls into 3 boxes and also 8 other balls into 2 other boxes, such that each box has a nonzero amount of balls.

 There are ${8+2-3 \choose 2}$ ways for the first part, and ${8+1-2 \choose 1}$ ways for the second part, by stars and bars.

 The answer is $2\cdot {7 \choose 2} \cdot {7 \choose 1} = \boxed{294}$.

 Feel free to edit this solution

 





\textbf{Solution 4}

Starting at the origin, you can either first go up or to the right. If you go up first, you will end on the side opposite to it (the right side) and if you go right first, you will end up on the top. It can then be observed that if you choose the turning points in the middle $7 \times 7$ grid, that will automatically determine your start and ending points. For example, in the diagram if you choose the point $(3,2)$ and $(5,3)$, you must first move three up or two right, determining your first point, and move 5 up or 3 right, determining your final point. Knowing this is helpful because if we first move anywhere horizontally, we have $7$ points on each column to choose from and starting from left to right, we have $6,5,4,3,2,1$ points on that row to choose from. This gives us $7(6)+7(5)+7(4)+7(3)+7(2)+7(1)$ which simplifies to $7\cdot21$. The vertical case is symmetrical so we have $7\cdot21\cdot2 = \boxed{294}$

 





\textbf{Solution 5}

As in Solution 1, there are two cases: $RURUR$ or $URURU$. We will work with the first case and multiply by $2$ at the end. We use stars and bars; we can treat the $R$s as the stars and the $U$s as the bars. However, we must also use stars and bars on the $U$s to see how many different patterns of bars we can create for the reds. We must have $1$ bar in $8$ blacks, so we use stars and bars on the equation 
 \[x + y = 8\]. However, each divider must have at least one black in it, so we do the change of variable $x' = x-1$ and $y' = x-1$. Our equation becomes 
 \[x' + y' = 6\]. By stars and bars, this equation has $\binom{6 + 2 - 1}{1} = 7$ valid solutions. Now, we use stars and bars on the reds. We must distribute two bars amongst the reds, so we apply stars and bars to 
 \[x + y + z = 8\]. Since each group must have one red, we again do a change of variables with $x' = x-1$, $y' = y-1$, and $z' = z-1$. We are now working on the equation 
 \[x' + y' + z' = 5\]. By stars and bars, this has $\binom{5 + 3 - 1}{2} = 21$ solutions. The number of valid paths in this case is the number of ways to create the bars times the number of valid arrangements of the stars given fixed bars, which equals $21 \cdot 7 = 147$. We must multiply by two to account for both cases, so our final answer is $147 \cdot 2 = \boxed{294}$.

 





\textbf{Solution 6}

We note that we either start with $R$ or $U$ and that these cases are symmetric, so we will simply consider the $R$ case and multiply by $2$ at the end. We make the key observation that making a path is the same as choosing two vertical lines and one horizontal line to mark our path; this gives us a unique path. Because we change directions exactly 4 times, we must travel at least 1 unit before turning, and we also must end by traveling at least one unit, so we cannot choose vertical lines that are on the edge of our grid. Similarly, we cannot do that for the horizontal line. Our final answer is simply choosing 2 vertical lines from 7 total, and choosing one horizontal line from 7 total, all multiplied by two. $2 \cdot \binom{7}{2} \cdot \binom{7}{1} = \boxed{294}.$

 





\textbf{Solution 7}


                    \begin{figure}[H]
                        \centering
                        \includegraphics[width=0.5\textwidth]{Solution/AIME_2024_I_6_Solution_Figure_1.jpeg}
                    \end{figure}
                    

 
Our path must either go right, up, right, up, and right (RURUR) or up, right, up, right, and up (URURU). We will consider the RURUR case and multiply by 2 since they're symmetrical.

 Label the $x$-axis from 0 to 8, starting at the origin and ending at the edge. For the first move (right), we can stop at point 1, 2, 3, 4, 5, 6, or 7.

 Let's consider the choices for the first and third rightward moves. If we stop at point 7 on the first rightward move, when we get to the third move, we will have 0 choices for choosing a point to go right until; no matter what we do, if we stop at point 7 on the first move, the path will never have more than three turns (see diagram). If we stop at point 6, we will have 1 choice for the third move (we can stop going right until one point before the right edge of the grid for the third move and our path will have 4 turns). If we stop at point 5, we can stop going right until one or two points before the right edge of the grid for the third move. As we continue to move down points, the number of choices for the third rightward move increases by 1. This pattern continues until we stop at point 1 on the first move and have 6 choices for the third move. So, for the first and third moves, we have a total of 
 \[0 + 1 + 2 + 3 + 4 + 5 + 6 = \frac{6 \cdot 7}{2} = 21\] choices.

 For the second move (upward), we always have 7 choices to stop at (the top, one point before the top, two points before the top... 1 point from the bottom), and the fourth move will always have 1 choice, which is the move that reaches the final point.

 Finally, we multiply by 2 to account for the URURU case, which gives us a total of $21 \cdot 7 \cdot 2 = \boxed{294}$  paths.

 \textbf{Fast Video Solution (no stars and bars)} 




	\item (\textit{AIME I 2024}) Find the largest possible real part of 
 \[(75+117i)z+\frac{96+144i}{z}\]where $z$ is a complex number with $|z|=4$. Here $i = \sqrt{-1}$.


 
 \textbf{Answer: }540



\textbf{Solution 1}

Let $z=4e^{i\theta}$. Then the expression becomes 
 \[(300+468i)e^{i\theta}+(24+36i)e^{-i\theta}\] The real part of this comes out to be $(300\cos \theta - 468 \sin \theta)+(24\cos \theta + 36 \sin \theta) = (324\cos \theta - 432 \sin \theta)$. Using Cauchy-Schwarz to maximize this, 
 \[(324\cos \theta - 432 \sin \theta)^2 \leq (\cos ^2 \theta+\sin ^2 \theta)(324^2+432^2) = (1)(324^2+432^2) = 540^2\] This value is attainable. Thus, the maximum is $\boxed{540}$.

 





\textbf{Solution 2a (Cauchy-Schwarz and vector algebra)}

Simplify rectangular form as in Solution 1 until we get $\text{Re}(w)=81a-108b = 27(3a-4b)$. 

 By Cauchy-Schwarz, to maximize $\text{Re}(w)$, the vector $z=[a,b]$ ( $|z| =4$) is $\frac{4}{|[3,-4]|}[3,-4]$. 

 We don't need to bash the arithmetic next, because the unit vector $u$  that maximizes $v \cdot u$ is $u=v/|v|$, so $v \cdot u= v\cdot v = |v|^2/|v| = |v|$, which here is just $\sqrt{3^2+(-4)^2} =5$.

 Combining what remains, we get answer $= 27 |z| |v|  = 27(4)(5)=\boxed{540}$.

 





\textbf{Solution 2b (Simple Analytic Geometry)}

Simplify rectangular form as in Solution 1 until we get $\text{Re}(w)=81x-108y = 27(3x-4y)$. 

 We also know $|z|=4$ or $x^2+y^2=16$. 

 By AM-GM or Cauchy-Schwarz, b = 4a/3, so 

 You can also prove this like so:We want to find the line $81x-108y=k$ tangent to circle $x^2+y^2=16$, which is perpendicular to the line connecting tangent point to circle's center $(0,0)$.

 Using the formula for (perpendicular) distance from a point to a line: $\frac{|ax+by+c|}{\sqrt{a^2+b^2}}=r$ we can substitute and get$\frac{|81(0)-108(0)-k|}{\sqrt{81^2+108^2}}=4$
 \begin{align*} \frac{k}{27\sqrt{3^2+4^2}}&=4 \\\frac{k}{135}&=4 \\k&=\boxed{540} \end{align*}


 





\textbf{Solution 2c (Dot product)}

Let $z = a + bi.$ Simplify until we get to maximizing $81a - 108b$ given $a^2 +b^2 = 16.$ We can write $81a - 108b$ as the \href{https://artofproblemsolving.com/wiki/index.php/Vector?srsltid=AfmBOoqdT29SM62x\_yUNcVA7yUtSecmXlpmyOtury4Rl4Y\_Bo-6fLXh4}{Dot Product} of two vectors:  
 \[81a - 108b = \langle 81 , -108 \rangle \cdot \langle a, b \rangle.\] From this, we have the line $b = -\frac{108}{81}a =  -\frac{4}{3}a$ and the circle $a^2 + b^2 = 16.$ We need to find their intersection $(a,b)$ such that $a$ and $b$ are maximized. The line intersects the circle at two points, but to maximize $81a - 108b$ we want $b$ to be negative and $a$ to be positive. So the point that works has coordinates $a = 2.4$ and $b = -3.2$. We get $81\cdot 2.4 + 108 \cdot 3.2 = \boxed{540}.$

 

                    \begin{figure}[H]
                        \centering
                        \includegraphics[width=0.5\textwidth]{Solution/AIME_2024_I_7_Solution_Figure_1.png}
                    \end{figure}
                    

 





\textbf{Solution 3}

Follow Solution 1 to get $81a-108b$. We can let $a=4\cos\theta$ and $b=4\sin\theta$ as $|z|=4$, and thus we have $324\cos\theta-432\sin\theta$. Furthermore, we can ignore the negative sign in front of the second term as we are dealing with sine and cosine, so we finally wish to maximize $324\cos\theta+432\sin\theta$ for obviously positive $\cos\theta$ and $\sin\theta$.

 
Using the previous fact, we can use the \href{https://artofproblemsolving.com/wiki/index.php?title=Cauchy-Schwarz\_Inequality}{Cauchy-Schwarz Inequality} to calculate the maximum. By the inequality, we have:

 $(324^2+432^2)(\cos^2\theta+\sin^2\theta)\ge(324\cos\theta+432\sin\theta)^2$

 $540^2\cdot1\ge(324\cos\theta+432\sin\theta)^2$

 $\boxed{540}\ge324\cos\theta+432\sin\theta$

 





\textbf{Solution 4 (Simple Quadratic Discriminant)}

Similar to the solutions above, we find that $Re((75+117i)z+\frac{96+144i}{z})=81a-108b=27(3a-4b)$, where $z=a+bi$. To maximize this expression, we must maximize $3a-4b$. Let this value be $x$. Solving for $a$ yields $a=\frac{x+4b}{3}$. From the given information we also know that $a^2+b^2=16$. Substituting $a$ in terms of $x$ and $b$ gives us $\frac{x^2+8bx+16b^2}{9}+b^2=16$. Combining fractions, multiplying, and rearranging, gives $25b^2+8xb+(x^2-144)=0$. This is useful because we want the maximum value of $x$ such that this quadratic has real roots which is easy to find using the discriminant. For the roots to be real, $(8x)^2-4(25)(x^2-144) \ge 0$. Now all that is left to do is to solve this inequality. Simplifying this expression, we get $-36x^2+14400 \ge 0$ which means $x^2 \le 400$ and $x \le 20$. Therefore the maximum value of $x$ is $20$ and $27 \cdot 20 = \boxed{540}$

 





\textbf{Solution 5 ("Completing the Triangle")}

First, recognize the relationship between the reciprocal of a complex number $z$ with its conjugate $\overline{z}$, namely:

 
 \[\frac{1}{z} \cdot \frac{\overline{z}}{\overline{z}} = \frac{\overline{z}}{|z|^2} = \frac{\overline{z}}{16}\]

 Then, let $z = 4(\cos\theta + i\sin\theta)$ and $\overline{z} = 4(\cos\theta - i\sin\theta)$.

 
 \begin{align*} Re \left ((75+117i)z+\frac{96+144i}{z} \right) &= Re\left ( (75+117i)z + (6+9i)\overline{z}    \right ) \\                                                &= 4 \cdot Re\left ( (75+117i)(\cos\theta + i\sin\theta) + (6+9i)(\cos\theta - i\sin\theta)    \right ) \\                                                &= 4 \cdot (75\cos\theta - 117\sin\theta + 6\cos\theta + 9\sin\theta) \\                                                &= 4 \cdot (81\cos\theta - 108\sin\theta) \\                                                &= 4\cdot 27 \cdot (3\cos\theta - 4\sin\theta) \end{align*}


 Now, recognizing the 3 and 4 coefficients hinting at a 3-4-5 right triangle, we "complete the triangle" by rewriting our desired answer in terms of an angle of that triangle $\phi$ where $\cos\phi = \frac{3}{5}$ and $\sin\phi = \frac{4}{5}$

 
 \begin{align*} 4\cdot 27 \cdot(3\cos\theta - 4\sin\theta) &= 4\cdot 27 \cdot 5 \cdot (\frac{3}{5}\cos\theta - \frac{4}{5}\sin\theta) \\                                                &= 540 \cdot (\cos\phi\cos\theta - \sin\phi\sin\theta) \\                                                &= 540 \cos(\theta + \phi) \end{align*}


 Since the simple trig ratio is bounded above by 1, our answer is $\boxed{540}$

 





\textbf{Solution 6 (Cauchy-Schwarz Inequality ) (Fastest)}

Follow as solution 1 would to obtain $81a + 108\sqrt{16-a^2}.$

 By the Cauchy-Schwarz Inequality, we have

 
 \[(a^2 + (\sqrt{16-a^2})^2)(81^2 + 108^2) \geq (81a + 108\sqrt{16-a^2})^2,\]

 so 

 
 \[4^2 \cdot 9^2 \cdot 15^2 \geq (81a + 108\sqrt{16-a^2})^2\]

 and we obtain that $81a + 108\sqrt{16-a^2} \leq 4 \cdot 9 \cdot 15 = \boxed{540}.$

 
- \href{https://artofproblemsolving.com/wiki/index.php/User:Spectraldragon8}{spectraldragon8}

 





\textbf{Solution 7 (Geometry)}

Follow solution 2 to get that we want to find the line $81a-108b=k$ tangent to circle $a^2+b^2=16$. The line turns into $a=\frac{k}{81}+\frac{4b}{3}$	Connect the center of the circle to the tangency point and the y-intercept of the line. Let the tangency point be $A$, the y-intercept be $C$, and the center be $B$. Drop the perpendicular from $A$ to $BC$ and call it $D$. Let $AD=3x$, $DC=4x$. Then, $BD=\sqrt{AB^2-AD^2}=\sqrt{16-9x^2}$. By similar triangles, get that $\frac{BD}{AD}=\frac{AD}{DC}$, so $\frac{\sqrt{16-9x^2}}{3x}=\frac{3x}{4x}$. Solve this to get that $x=\frac{16}{15}$, so $BC=\frac{20}{3}$ and $\frac{k}{81}=\frac{20}{3}$, so $k=\boxed{540}$~ryanbear

 





\textbf{Solution 8 (Euler's Formula and Trig Optimization)}

Because $|z|=4$, we can let $z=4e^{i\theta}$. Then, substituting $i=e^{\frac{i\pi}{2}}$, we get that the complex number is 

 \begin{align*}w&=4e^{i\theta}(75+117e^{\frac{i\pi}{2}})+\dfrac{1}{4}e^{-i\theta}(96+144e^{\frac{i\pi}{2}})\\&=300e^{i\theta}+468e^{i(\frac{\pi}{2}+\theta)}+24e^{-i\theta}+36e^{i(\frac{\pi}{2}-\theta)}\\\end{align*}


 We know that the $\text{Re}(e^{i\alpha})=\cos(\alpha)$ from Euler's formula, so applying this and then applying trig identities yields

 \begin{align*}\text{Re}(w)&=300\cos{(\theta)}+468\cos{(\dfrac{\pi}{2}+\theta)}+24\cos{(-\theta)}+36\cos{(\dfrac{\pi}{2}-\theta)}\\&=300\cos{(\theta)}-468\sin{(\theta)}+24\cos{(\theta)}+36\sin{(\theta)}\\&=324\cos{(\theta)}-432\sin{(\theta)}\\\implies \dfrac{1}{108}\text{Re}(w)&=3\cos{(\theta)}-4\sin{(\theta)}\\\end{align*}


 We can see that the right-hand side looks an awful lot like the sum of angles formula for cosine, but 3 and 4 don't satisfy the pythagorean identity. To make them do so, we can divide everything by $\sqrt{3^2+4^2}=5$ and set $\cos{(\alpha)}::=\frac{3}{5}$ and $\sin{(\alpha)}::=\frac{4}{5}$. Now we have that 
 \[\dfrac{1}{540}\text{Re}(w)=\cos{(\theta+\alpha)}\]Obviously the maximum value of the right hand side is 1, so the maximum value of the real part is $\boxed{540}$.

 





\textbf{Solution 9 (Calc semi-bash)}

Let $c$ denote value of the above expression such that $\mathsf{Re} (c)$ is maximized. We write $z=4e^{i\theta}$ and multiply the second term in the expression by $\overline{z} = 4e^{-i\theta},$ turning the expression into
 \[4e^{i\theta}(75+117i) + \frac{(96 + 144i)\cdot 4e^{-i\theta}}{4e^{i\theta}\cdot 4e^{-i\theta}} = 300e^{i\theta} + 468ie^{i\theta} + (24+ 36i)e^{-i\theta}.\]Now, we write $e^{i\theta} = \cos\theta + i\sin\theta$. Since $\cos$ is even and $\sin$ is odd,
 \begin{align*} &300(\cos\theta + i\sin\theta) +468i + (24+36i)(\cos\theta -i\sin\theta) \\ \iff & \mathsf{Re}(c) = 324\cos\theta -468\sin\theta \end{align*}
We want to maximize this expression, so we take its derivative and set it equal to $0$ (and quickly check the second derivative for inflection points):
 \begin{align*}     &\mathsf{Re}(c) = 108\left(3\cos\theta - 4\sin\theta\right)\\     \frac{d}{d\theta} &\mathsf{Re}(c) = -324\sin\theta -468\cos\theta = 0, \end{align*}
so $\tan\theta = -\dfrac{468}{324} = -\dfrac{4}{3},$ which is reminiscent of a $3-4-5$ right triangle in the fourth quadrant (side lengths of $3, -4, 5$). Since $\tan\theta = -\frac{4}{3},$ we quickly see that $\sin\theta = -\dfrac{4}{5}$ and $\cos\theta = \dfrac{3}{5}.$ Therefore,
 \begin{align*}  \mathsf{Re}(c) &= 108\left(3\cos\theta - 4\sin\theta \right) =  108\left(\frac{9}{5} + \frac{16}{5} \right)  = 108\cdot 5 = \boxed{\textbf{(540)}} \end{align*}


 -Benedict T (countmath1)

 





\textbf{Solution 10 (Lagrange multipliers)}

With $z = a + bi$ such that $a^2 + b^2 = 16,$ we have 
 \begin{align*} (75 +117i)(a + bi) + \frac{48}{a + bi}  (2 + 3i) &= 75a + 75bi + 117ai - 117b + \frac{48}{a + bi}(2 + 3i) \\ &= 75a - 117b + (117a + 75b)i + 48 (2 + 3i) \cdot \frac{a - bi}{16} \\ &= 75a - 117b + (117a + 75b)i + 3 (2 + 3i)(a - bi) \end{align*}
where we use $z^{-1} = \frac{\bar z}{|z|^2}.$ With $3 (2 + 3i)(a - bi) = 3 [2a - 2bi + 3ai + 3b] = 6a +9b +9ai-6bi,$ the expression becomes $81a-108b+ (126a + 69b)i$ and we would like to maximize $81a - 108b = 9(9a - 12b) = 27(3a - 4b)$ with $a^2 + b^2 = 16.$With $f(a, b) = 3a - 4b$ and $g(a, b) = a^2 + b^2 = 16,$ we have 
 \begin{align*} 3 = 2\lambda a, \quad -4 = 2\lambda b \implies -\frac{3}{4} = \frac ab \implies -3b = 4a \implies b = -\frac 43 a\end{align*}
so 
 \[a^2 + \frac{16}{9}a^2 = \frac{25}{9}a^2 = 16 \implies \frac{5}{3}a = 4 \implies a = \frac {12}5, b = -\frac{16}{5}\]and we have $3a - 4b = \frac{36}{5} + \frac{64}{5} = 20,$ so the maximum is $27 \cdot 20 = \boxed{540}.$-centslordm

 





\textbf{Solution 11 (basically Lagrange but easier-ish)}

Proceeding as with Solution 10, we aim to maximize $27(3a-4b)$ under the constraint $a^2+b^2-16=0$. It is a well-known result of Lagrange multipliers that a linear function is maximized under a circle when the values of the variables are proportional to their coefficients; that is, in our case, $a=3x$ and $b=-4x$. (Technically $a=27\cdot3x$ and $b=27(-4)x$, but it's easier to use $a=3x$ and $b=-4x$.) 

 Then, $(3x)^2+(-4x)^2=25x^2=16$, so $x=\dfrac45$ and we have $a=\dfrac{12}5$ and $b=-\dfrac{16}5$. This yields 

 
 \[27\left(3\left(\dfrac{12}5\right)-4\left(-\dfrac{16}5\right)\right)=27\left(\dfrac{36}5+\dfrac{64}5\right)=27\cdot\dfrac{100}5=\boxed{540}~.\] 

 QED. $\Box$

 





\textbf{Solution 12 (Wave function)}

Note that we can scale down the expression by a factor of $3$, namely, we compute the maximum possible real part of ${25+39i}z+{32+48i\over z}$ for $|z|=4$ and multiply this result by 3. Write $z=4(\cos t+i\sin t)$ we have

 \begin{align}(25+39i)z + \frac{32+48i}{z}&= 4(25+39i)(\cos t+i\sin t) + \frac{32+48i}{4(\cos t+i\sin t)} \\&= (100+156i)(\cos t+i\sin t) + (8+12i)(\cos t-i\sin t)\end{align}

 The real part of $(100+156i)(\cos t+i\sin t)+(8+12i)(\cos t-i\sin t)$ is then 
 \[108\cos t-144\sin t=36(3\cos t-4\sin t)\]

 The function $3\cos t-4\sin t$ oscillates with amplitude $5$, so the maximum value of the scaled-down expression is $36\cdot 5=180$. Hence, our requested answer is $180\cdot 3=\boxed{540}$.

 





\textbf{Solution 13 (Conjugate)(simplest)}

Suppose $z = 4e^{i\theta}$

 \begin{align}& \phantom{\text{ = }}\mathrm{Re}\left((75 + 117i)z + \dfrac{96 + 144i}{z}\right)\\& = \mathrm{Re}\left((300 + 468i)e^{i\theta} + (24 + 36i)e^{-i\theta}\right)\\& = \mathrm{Re}\left((300 + 468i)e^{i\theta} + (24 - 36i)e^{i\theta}\right)\\& = \mathrm{Re}\left((324 + 432i)e^{i\theta}\right)\\& \leq \sqrt{324^2 + 432^2}\\& = \boxed{540}\end{align}

 \textbf{Remark 13.1} 
                \begin{itemize}
                    \item $\overline{z_1z_2} = \overline{z_1} \cdot \overline{z_2}$,
                \end{itemize}


                
                \begin{itemize}
                    \item $24 + 36i, 24 - 36i$ and $e^{-i\theta}, e^{i\theta}$ are conjugates,
                \end{itemize}


                
                \begin{itemize}
                    \item The real parts of conjugates are equal,
                \end{itemize}


                
                \begin{itemize}
                    \item $\mathrm{Re}(z) \leq |z|$.
                \end{itemize}


                \textbf{Remark 13.2} The idea comes from the geometric meaning of this problem. The only difficulty is that the rotation directions are opposite. To make the rotations same, using conjugate is a good idea without changing the real part. Finally we can get a rotating parallelogram as shown in the following graph. Just find the length of the diagonal.

 
                    \begin{figure}[H]
                        \centering
                        \includegraphics[width=0.5\textwidth]{Solution/AIME_2024_I_7_Solution_Figure_2.png}
                    \end{figure}
                    

 





\textbf{Solution 14 (Dot Product Maximality, The Hadamard Vector)}

Let $z=4e^{i \theta}$. We have

 
 \[(75+117i)(4e^{i \theta}) + (96+144i)(\frac{1}{4}e^{-i\theta})\]
 \[(300+468i)(\cos(\theta) + i\sin(\theta)) + (24+36i)(\cos(\theta) - i\sin(\theta))\]
 \[(300\cos(\theta)-468\sin(\theta))+i(468\cos(\theta)+300\sin(\theta))+(24\cos(\theta)+36\sin(\theta)) +i(36\cos(\theta)-24\sin(\theta))\]
 \[(324\cos(\theta)-432\sin(\theta))+i(504\cos(\theta)+276\sin(\theta))\]

 From here, we need to maximize

 
 \[324\cos(\theta) - 432\sin(\theta)\].

 Recall for two vectors $\textbf{a} = (a_1,a_2)$ and $\textbf{b} = (b_1,b_2)$, the \textbf{dot product property} states $\textbf{ab} = (a_1, a_2) \cdot (b_1,b_2) = a_1b_1 + a_2b_2$. 

 Rewriting $324\cos(\theta) - 432\sin(\theta)$ as a dot product, we have $(324, -432) \cdot (\cos(\theta), \sin(\theta))$, so one vector yields $\textbf{u} = (324, -432)$ and $\textbf{v} = (\cos(\theta), \sin(\theta))$. To find the maximum $| \textbf{uv} |$, we require the maximalization of the product to form a vector of maximal magnitude. Consider the Hadamard product, the resulting vector formed via the dot product. Then, this vector is maximal when both vectors face the same direction. The largest possible vector direction for $\textbf{v}$ always yields a magnitude 1. For the maximality of $\textbf{u}$, consider $|\textbf{u}|  = \sqrt{324^2 + 432^2}$. Divide both by 108 and multiply at the end, obtaining $108 \cdot \sqrt{3^2 + 4^2} = 108 \cdot 5 = 540$. So, the max value of $| \textbf{uv} |$ is $540 \cdot 1 = \boxed{540}$.

 




	\item (\textit{AIME I 2024}) Eight circles of radius $34$ can be placed tangent to $\overline{BC}$ of $\triangle ABC$ so that the circles are sequentially tangent to each other, with the first circle being tangent to $\overline{AB}$ and the last circle being tangent to $\overline{AC}$, as shown. Similarly, $2024$ circles of radius $1$ can be placed tangent to $\overline{BC}$ in the same manner. The inradius of $\triangle ABC$ can be expressed as $\frac{m}{n}$, where $m$ and $n$ are relatively prime positive integers. Find $m+n$.


 
                    \begin{figure}[H]
                        \centering
                        \includegraphics[width=0.5\textwidth]{Problem/AIME_2024_I_8_Problem_Figure_1.png}
                    \end{figure}
                    


 
 \textbf{Answer: }197



\textbf{Solution 1}

Draw an altitude from both end circles of the diagram with the circles of radius one, and call the lengths at the end cut off by the altitudes the altitudes of the circles down to $BC$ $a$ and $b$. Now we have the length of side $BC$ of being $(2)(2022)+1+1+a+b$. However, the side $BC$ can also be written as $(6)(68)+34+34+34a+34b$, due to similar triangles from the second diagram. If we set the equations equal, we have $\frac{1190}{11} = a+b$. Call the radius of the incircle $r$, then we have the side BC to be $r(a+b)$. We find $r$ as $\frac{4046+\frac{1190}{11}}{\frac{1190}{11}}$, which simplifies to $\frac{10+((34)(11))}{10}$,so we have $\frac{192}{5}$, which sums to $\boxed{197}$.

 





\textbf{Solution 2}

Assume that $ABC$ is isosceles with $AB=AC$.

 If we let $P_1$ be the intersection of $BC$ and the leftmost of the eight circles of radius $34$, $N_1$ the center of the leftmost circle, and $M_1$ the intersection of the leftmost circle and $AB$, and we do the same for the $2024$ circles of radius $1$, naming the points $P_2$, $N_2$, and $M_2$, respectively, then we see that $BP_1N_1M_1\sim BP_2N_2M_2$. The same goes for vertex $C$, and the corresponding quadrilaterals are congruent.

 Let $x=BP_2$. We see that $BP_1=34x$ by similarity ratios (due to the radii). The corresponding figures on vertex $C$ are also these values. If we combine the distances of the figures, we see that $BC=2x+4046$ and $BC=68x+476$, and solving this system, we find that $x=\frac{595}{11}$.

 If we consider that the incircle of $\triangle ABC$ is essentially the case of $1$ circle with $r$ radius (the inradius of $\triangle ABC$), we can find that $BC=2rx$. From $BC=2x+4046$, we have:

 
 \[r=1+\frac{2023}{x}=1+\frac{11\cdot2023}{595}=1+\frac{187}{5}=\frac{192}{5}\]

 Thus the answer is $192+5=\boxed{197}$.

 





\textbf{Solution 3}

Let $x = \cot{\frac{B}{2}} + \cot{\frac{C}{2}}$. By representing $BC$ in two ways, we have the following:
 \[34x + 7\cdot 34\cdot 2 = BC\]
 \[x + 2023 \cdot 2 = BC\]

 Solving we find $x = \frac{1190}{11}$. Now draw the inradius, let it be $r$. We find that $rx =BC$, hence 
 \[xr = x + 4046 \implies r-1 = \frac{11}{1190}\cdot 4046 = \frac{187}{5}.\]Thus 
 \[r = \frac{192}{5} \implies \boxed{197}.\]~AtharvNaphade

 





\textbf{Solution 4}

First, let the circle tangent to $AB$ and $BC$ be $O$ and the other circle that is tangent to $AC$ and $BC$ be $R$. Let $x$ be the distance from the tangency point on line segment $BC$ of the circle $O$ to $B$. Also, let $y$ be the distance of the tangency point of circle $R$ on the line segment $BC$ to point $C$. Realize that we can let $n$ be the number of circles tangent to line segment $BC$ and $r$ be the corresponding radius of each of the circles. Also, the circles that are tangent to $BC$ are similar. So, we can build the equation $BC = (x+y+2(n-1)) \times r$. Looking at the given information, we see that when $n=8$, $r=34$, and when $n=2024$, $r=1$, and we also want to find the radius $r$ in the case where $n=1$. Using these facts, we can write the following equations:

 $BC = (x+y+2(8-1)) \times 34 = (x+y+2(2024-1)) \times 1 = (x+y+2(1-1)) \times r$

 We can find that $x+y = \frac{1190}{11}$ . Now, let $(x+y+2(2024-1)) \times 1 = (x+y+2(1-1)) \times r$. 

 Substituting $x+y = \frac{1190}{11}$ in, we find that 
 \[r = \frac{192}{5} \implies \boxed{197}.\]

 





\textbf{Solution 5 (one variable)}

Define $I, x_1, x_8, y_1, y_{2024}$ to be the incenter and centers of the first and last circles of the $8$ and $2024$ tangent circles to $BC,$ and define $r$ to be the inradius of triangle $\bigtriangleup ABC.$ We calculate $\overline{x_1x_8} = 34 \cdot 14$ and $\overline{y_1y_{2024}} = 1 \cdot 4046$ because connecting the center of the circles voids two extra radii.

 
We can easily see that $B, x_1, y_1,$ and $I$ are collinear, and the same follows for $C, x_8, y_{2024},$ and $I$ (think angle bisectors).

 
We observe that triangles $\bigtriangleup I x_1 x_8$ and $\bigtriangleup I y_1 y_{2024}$ are similar, and therefore the ratio of the altitude to the base is the same, so we note

 
 \[\frac{\text{altitude}}{\text{base}} = \frac{r-34}{34\cdot 14} = \frac{r-1}{1\cdot 4046}.\]

 
Solving yields $r = \frac{192}{5},$ so the answer is $192+5 = \boxed{197}.$

 -\href{https://artofproblemsolving.com/wiki/index.php/User:Spectraldragon8}{spectraldragon8}

 




	\item (\textit{AIME I 2024}) Let $A$, $B$, $C$, $D$ be points on the hyperbola  $\frac{x^2}{20}-\frac{y^2}{24}=1$ such that $ABCD$ is a rhombus whose diagonals intersect at the origin. Find the largest number less than $BD^2$ for all rhombuses $ABCD$.


 
 \textbf{Answer: }480



\textbf{Solution 1}

A quadrilateral is a rhombus if and only if its two diagonals bisect each other and are perpendicular to each other. The first condition is automatically satisfied because of the hyperbola's symmetry about the origin. To satisfy the second condition, we set $BD$ as the line $y = mx$ and $AC$ as $y = -\frac{1}{m}x.$ Because the hyperbola has asymptotes of slopes $\pm \frac{\sqrt6}{\sqrt5},$ we have $m, -\frac{1}{m} \in \left(-\frac{\sqrt6}{\sqrt5}, \frac{\sqrt6}{\sqrt5}\right).$ This gives us $m^2 \in \left(\frac{5}{6}, \frac{6}{5}\right).$

 
Plugging $y = mx$ into the equation for the hyperbola yields $x^2 = \frac{120}{6-5m^2}$ and $y^2 = \frac{120m^2}{6-5m^2}.$ By symmetry of the hyperbola, we know that $\left(\frac{BD}{2}\right)^2 = x^2 + y^2,$ so we wish to find a lower bound for $x^2 + y^2 = 120\left(\frac{1+m^2}{6-5m^2}\right).$ This is equivalent to minimizing $\frac{1+m^2}{6-5m^2} = -\frac{1}{5} + \frac{11}{5(6-5m^2)}$. It's then easy to see that this expression increases with $m^2,$ so we plug in $m^2 = \frac{5}{6}$ to get $x^2+y^2 > 120,$ giving $BD^2 > \boxed{480}.$

 





\textbf{Solution 2}

Assume that $AC$ is the asymptote of the hyperbola, in which case $BD$ is minimized. The expression of $BD$ is $y=-\sqrt{\frac{5}{6}}x$. Thus, we could get $\frac{x^2}{20}-\frac{y^2}{24}=1\implies x^2=\frac{720}{11}$. The desired value is $4\cdot \frac{11}{6}x^2=480$. This case can't be achieved, so all $BD^2$ would be greater than $\boxed{480}$

 


 


 




	\item (\textit{AIME I 2024}) Let $\triangle ABC$ have side lengths $AB=5$, $BC=9$, $CA=10$. The tangents to circumcircle of $\triangle ABC$ at $B$ and $C$ intersect at point $D$, and $\overline{AD}$ intersects the circumcircle at $P \neq A$. The length of $AP$ is equal to $\frac{m}{n}$, where $m$ and $n$ are relatively prime integers. Find $m + n$.


 
 \textbf{Answer: }113



\textbf{Solution 1}

We have $\let\angle BCD = \let\angle CBD = \let\angle A$ from the tangency condition. With LoC we have $\cos(A) = \frac{25+100-81}{2*5*10} = \frac{11}{25}$ and $\cos(B) = \frac{81+25-100}{2*9*5} = \frac{1}{15}$. Then, $CD = \frac{\frac{9}{2}}{\cos(A)} = \frac{225}{22}$. Using LoC we can find $AD$: $AD^2 = AC^2 + CD^2 - 2(AC)(CD)\cos(A+C) = 10^2+(\frac{225}{22})^2 + 2(10)\frac{225}{22}\cos(B) = 100 + \frac{225^2}{22^2} + 2(10)\frac{225}{22}*\frac{1}{15} = \frac{5^4*13^2}{484}$. Thus, $AD = \frac{5^2*13}{22}$. By Power of a Point, $DP*AD = CD^2$ so $DP*\frac{5^2*13}{22} = (\frac{225}{22})^2$ which gives $DP = \frac{5^2*9^2}{13*22}$. Finally, we have $AP = AD - DP = \frac{5^2*13}{22} - \frac{5^2*9^2}{13*22} = \frac{100}{13}$. So the answer is $\boxed{113}$.

 





\textbf{Solution 2 (no knowledge of symmedian, reduced geo)}


                    \begin{figure}[H]
                        \centering
                        \includegraphics[width=0.5\textwidth]{Solution/AIME_2024_I_10_Solution_Figure_1.png}
                    \end{figure}
                    This solution is longer given the fact the reader is unfamiliar with Symmedians.

 Through the two tangents theorem $CD = DB$ and so $\angle DCB = \angle DBC$. Also, through Tangent Chord it follows that $\angle A  = \angle DCB = \angle DBC$. Hence, drop the altitude $D$ to base $CB$ and call it $M$. Conjoin median $AM$. 

 Conjoin $BP$ and $PC$ to form a cyclic quadrilateral $ABPC$. Firstly, notice that $\angle ACB = \angle APB$. Secondly, extend median $AM$ to meet the circumcircle $\omega$ at $Q$, and notice that $CQPB$ is an isosceles trapezoid, and therefore minor arcs $BP$ and $CQ$ have the same angle, and so their inscribed angles $CAQ$ and $PAB$ are equal, or in other terms, $CAM$ and $PAB$ are equal. Therefore, $\triangle{ABP}\sim \triangle{AMC}$. We can show the isosceles trapezoid quite easily. Firstly, we prove that the figure is a trapezoid. Note that $CPQ$ and $CBP$ are the same angle, and PB is a line that intersects through lines $QP$ and $CB$. Through Converse of Corresponding Angles Theorem, the two lines $QP$ and $CB$ are parallel, so the figure is a trapezoid by definition. We now show it is isosceles. We have major arc $BC$ inscribed by both $CPB$ and $CQB$, so $\angle CPB = \angle CQB$, and correspondingly $\angle QCB = \angle PBC$. Rest is as follows.

 By Apollonius's theorem, $AM=\frac{13}{2}$. Through similarity, $AP=\frac{100}{13}\implies \boxed{113}$

 





\textbf{Solution 3}


                    \begin{figure}[H]
                        \centering
                        \includegraphics[width=0.5\textwidth]{Solution/AIME_2024_I_10_Solution_Figure_2.png}
                    \end{figure}
                    We know $AP$ is the symmedian (see \href{https://artofproblemsolving.com/wiki/index.php?title=Symmedians,\_Lemoine\_point}{ Symmedian and tangents }) ,

 which implies that $\triangle{ABP}\sim \triangle{AMC}$ where $M$ is the midpoint of $BC$.

 By Apollonius' theorem, $AM=\frac{13}{2}$. 

 Thus, we have $\frac{AP}{AC}=\frac{AB}{AM}, AP=\frac{100}{13}\implies \boxed{113}$

 





\textbf{Solution 4}

Extend sides $\overline{AB}$ and $\overline{AC}$ to points $E$ and $F$, respectively, such that $B$ and $C$ are the feet of the altitudes in $\triangle AEF$. Denote the feet of the altitude from $A$ to $\overline{EF}$ as $X$, and let $H$ denote the orthocenter of $\triangle AEF$. Call $M$ the midpoint of segment $\overline{EF}$. By the Three Tangents Lemma, we have that $MB$ and $MC$ are both tangents to $(ABC)$ $\implies$ $M = D$, and since $M$ is the midpoint of $\overline{EF}$, $MF = MB$. Additionally, by angle chasing, we get that: 
 \[\angle ABC \cong \angle AHC \cong \angle EHX\]Also, 
 \[\angle EHX = 90 ^\circ - \angle HEF = 90 ^\circ - (90 ^\circ - \angle AFE) = \angle AFE\] Furthermore, 
 \[AB = AF \cdot \cos(A)\]From this, we see that $\triangle ABC \sim \triangle AFE$ with a scale factor of $\cos(A)$. By the Law of Cosines, 
 \[\cos(A) = \frac{10^2 + 5^2 - 9^2}{2 \cdot 10 \cdot 5} = \frac{11}{25}\] Thus, we can find that the side lengths of $\triangle AEF$ are $\frac{250}{11}, \frac{125}{11}, \frac{225}{11}$. Then, by Stewart's theorem, $AM = \frac{13 \cdot 25}{22}$. By Power of a Point, 
 \[\overline{MB} \cdot \overline{MB} = \overline{MA} \cdot \overline{MP}\]
 \[\frac{225}{22} \cdot \frac{225}{22} = \overline{MP} \cdot \frac{13 \cdot 25}{22} \implies \overline{MP} = \frac{225 \cdot 9}{22 \cdot 13}\]Thus, 
 \[AP = AM - MP = \frac{13 \cdot 25}{22} - \frac{225 \cdot 9}{22 \cdot 13} = \frac{100}{13}\]Therefore, the answer is $\boxed{113}$.

 





\textbf{Solution 5 (LoC spam)}

Connect lines $\overline{PB}$ and $\overline{PC}$. From the angle by tangent formula, we have $\angle PBD = \angle DAB$. Therefore by AA similarity, $\triangle PBD \sim \triangle BAD$. Let $\overline{BP} = x$. Using ratios, we have 
 \[\frac{x}{5}=\frac{BD}{AD}.\] Similarly, using angle by tangent, we have $\angle PCD = \angle DAC$, and by AA similarity, $\triangle CPD \sim \triangle ACD$. By ratios, we have 
 \[\frac{PC}{10}=\frac{CD}{AD}.\] However, because $\overline{BD}=\overline{CD}$, we have 
 \[\frac{x}{5}=\frac{PC}{10},\] so $\overline{PC}=2x.$ Now using Law of Cosines on $\angle BAC$ in triangle $\triangle ABC$, we have 
 \[9^2=5^2+10^2-100\cos(\angle BAC).\] Solving, we find $\cos(\angle BAC)=\frac{11}{25}$. Now we can solve for $x$. Using Law of Cosines on $\triangle BPC,$ we have \begin{align*}81&=x^2+4x^2-4x^2\cos(180-\angle BAC) \\ &= 5x^2+4x^2\cos(BAC). \\\end{align*}


 Solving, we get $x=\frac{45}{13}.$ Now we have a system of equations using Law of Cosines on $\triangle BPA$ and $\triangle CPA$, 
 \[AP^2=5^2+\left(\frac{45}{13}\right)^2 -(10) \left(\frac{45}{13} \right)\cos(ABP)\]
 \[AP^2=10^2+4 \left(\frac{45}{13} \right)^2 + (40) \left(\frac{45}{13} \right)\cos(ABP).\]

 Solving, we find $\overline{AP}=\frac{100}{13}$, so our desired answer is $100+13=\boxed{113}$.

 





\textbf{Solution 6}

Following from the law of cosines, we can easily get $\cos A = \frac{11}{25}$, $\cos B = \frac{1}{15}$, $\cos C = \frac{13}{15}$.

 Hence, $\sin A = \frac{6 \sqrt{14}}{25}$, $\cos 2C = \frac{113}{225}$, $\sin 2C = \frac{52 \sqrt{14}}{225}$.Thus, $\cos \left( A + 2C \right) = - \frac{5}{9}$.

 Denote by $R$ the circumradius of $\triangle ABC$.In $\triangle ABC$, following from the law of sines, we have $R = \frac{BC}{2 \sin A} = \frac{75}{4 \sqrt{14}}$.

 Because $BD$ and $CD$ are tangents to the circumcircle $ABC$, $\triangle OBD \cong \triangle OCD$ and $\angle OBD = 90^\circ$.Thus, $OD = \frac{OB}{\cos \angle BOD} = \frac{R}{\cos A}$.

 In $\triangle AOD$, we have $OA = R$ and $\angle AOD = \angle BOD + \angle AOB = A + 2C$.Thus, following from the law of cosines, we have

 \begin{align*}AD & = \sqrt{OA^2 + OD^2 - 2 OA \cdot OD \cos \angle AOD} \\& = \frac{26 \sqrt{14}}{33} R.\end{align*}


 
Following from the law of cosines,

 \begin{align*}\cos \angle OAD & = \frac{AD^2 + OA^2 - OD^2}{2 AD \cdot OA} \\& = \frac{8 \sqrt{14}}{39} .\end{align*}


 
Therefore,

 \begin{align*}AP & = 2 OA \cos \angle OAD \\& = \frac{100}{13} .\end{align*}


 
Therefore, the answer is $100 + 13 = \boxed{\textbf{(113) }}$.

 





\textbf{Solution 7}

Note that since P is a symmedian, $(AP;BC)$ are harmonic. As a result, $\frac{BP}{PC} = \frac{BA}{CA}$. As a result, $2(BP) = PC$. Call $BP = x$. Then, $PC = 2x$. Since $cos A = \frac{11}{25}$, $cos BPC = - \frac{11}{25}$. Use LOC to find $x = \frac{45}{13}$. Finish with Ptolemy on ABPC, and finish to get $\frac{100}{13}$.

 





\textbf{Solution 8 (1 min solve)}

Note that $AD$ is the A-symmedian in $\triangle ABC$. Let $M$ be the midpoint of $BC$. It is a well known property that $\triangle ABP \sim \triangle AMC$. Therefore, using the median length formula, 
 \[AM = \frac{\sqrt{200+50-81}}{2} = \frac{13}{2},\]so 
 \[\frac{\frac{13}{2}}{10} = \frac{5}{AP} \implies AP = \frac{100}{13},\]so our answer is $\boxed{113}$.

 





\textbf{Solution 9 (Inversion)}

Take a $\sqrt{bc}$ inversion, and note that the circumcircle of $\triangle ABC$ gets mapped to $\overline{BC}$. Therefore, $P$ is mapped to a point on $\overline{BC}$. Also, the tangents map to circles that pass through $A$ and are tangent to $\overline{BC}$ at $B$ or $C$. The radical axis of the two circles therefore passes through the midpoint of $\overline{BC}$, and thus $P$ is mapped to the midpoint of $\overline{BC}$. By Stewart's, $AM=\frac{13}{2}$, so $AP=\frac{AB*AC}{AM}=\frac{100}{13}$. Therefore, the answer is $\boxed{113}$. 

 





\textbf{Solution 10 (Coordinate Geometry)}

Let the circumcircle of $\Delta ABC$ be $\omega$. By the Law of Cosines and Law of Sines, we have
 \[\cos{A}=\frac{11}{25},\sin{A}=\frac{6\sqrt{14}}{25}, R=\frac{75}{4\sqrt{14}}.\]Let $M$ be the midpoint of $BC$. We see that $\angle A=\angle DOB=\angle MBD$. It follows that $OM=\frac{33\sqrt{14}}{56},MD=\frac{27\sqrt{14}}{11}$. Let point $H$ be the foot of the perpendicular from point $A$ to segment $BC$. By the Pythagorean theorem, $MH=\frac{25}{6}$. We establish a coordinate system with $M$ being the origin. Here is our setup:\begin{align*}A&=\left(\frac{25}{6},\frac{4\sqrt{14}}{3}\right)\\B&=\left(\frac{9}{2},0\right)\\D&=\left(0,-\frac{27\sqrt{14}}{11}\right)\\M&=(0,0)\\P&=(x_1,y_1)\\AD&:y=\frac{10\sqrt{14}}{11}x-\frac{27\sqrt{14}}{11}\\\omega&:x^2+\left(y-\frac{33\sqrt{14}}{56}\right)^2=\frac{5625}{224}\end{align*}
From the equation of $\omega$ and $AD$, we have
 \[x^2+\left(\frac{10\sqrt{14}}{11}x-\frac{1875\sqrt{14}}{616}\right)^2=\frac{5625}{224}.\]By Vieta's formula, it follows that
 \[x_1+\frac{25}{6}=\frac{9375}{1521}.\]Note that the length of $AP$ can be computed by multiplying the secant value of $\theta$, the angle formed between line $AP$ and the $x$-axis, and the change in $x$-coordinates between $A$ and $P$. We see that\begin{align*}\Delta x=\left\lvert x_1-\frac{25}{6}\right\rvert=\frac{25}{3}-\frac{9375}{1521}=\frac{1100}{507}\\\tan{\theta}=\frac{10\sqrt{14}}{11}\implies\cos{\theta}=\frac{11}{39}\end{align*}
Therefore, 
 \[AP=\frac{1100}{507}\cdot\frac{39}{11}=\frac{100}{13}.\]The requested sum is $100+13=\boxed{113}$.

 




	\item (\textit{AIME I 2024}) Each vertex of a regular octagon is independently colored either red or blue with equal probability. The probability that the octagon can then be rotated so that all of the blue vertices end up at positions where there had been red vertices is $\frac{m}{n}$, where $m$ and $n$ are relatively prime positive integers. Find $m+n$.


 
 \textbf{Answer: }371



\textbf{Solution 1}

Notice that the question's condition mandates all blues to go to reds, but reds do not necessarily have to go to blue. Let us do casework on how many blues there are. 

 If there are no blues whatsoever, there is only one case. This case is valid, as all of the (zero) blues have gone to reds. (One could also view it as: the location of all the blues now were not previously red.) Thus, we have $1$. 

 If there is a single blue somewhere, there are $8$ cases - where can the blue be? Each of these is valid. 

 If there are two blues, again, every case is valid, and there are $\dbinom82=28$ cases. 

 If there are three blues, every case is again valid; there are $\dbinom83=56$ such cases. 

 The case with four blues is trickier. Let us look at all possible subcases. 

 If all four are adjacent (as in the diagram below), it is obvious: we can simply reverse the diagram (rotate it by $4$ units) to achieve the problem's condition. There are $8$ possible ways to have $4$ adjacent blues, so this subcase contributes $8$. 
                    \begin{figure}[H]
                        \centering
                        \includegraphics[width=0.5\textwidth]{Solution/AIME_2024_I_11_Solution_Figure_1.png}
                    \end{figure}
                    

 If three are adjacent and one is one away (as shown in the diagram below), we can not rotate the diagram to satisfy the question. This subcase does not work. 

 
                    \begin{figure}[H]
                        \centering
                        \includegraphics[width=0.5\textwidth]{Solution/AIME_2024_I_11_Solution_Figure_2.png}
                    \end{figure}
                    

 If three are adjacent and one is two away, obviously it is not possible as there is nowhere for the three adjacent blues to go. 

 
                    \begin{figure}[H]
                        \centering
                        \includegraphics[width=0.5\textwidth]{Solution/AIME_2024_I_11_Solution_Figure_3.png}
                    \end{figure}
                    

 If there are two adjacent pairs that are $1$ apart, it is not possible since we do not have anywhere to put the two pairs. 

 
                    \begin{figure}[H]
                        \centering
                        \includegraphics[width=0.5\textwidth]{Solution/AIME_2024_I_11_Solution_Figure_4.png}
                    \end{figure}
                    

 If there are two adjacent pairs that are $2$ apart, all of these cases are possible as we can rotate the diagram by $2$ vertices to work. There are $4$ of these cases. 

 
                    \begin{figure}[H]
                        \centering
                        \includegraphics[width=0.5\textwidth]{Solution/AIME_2024_I_11_Solution_Figure_5.png}
                    \end{figure}
                    

 If there is one adjacent pair and there are two separate ones each a distance of $1$ from the other, this case does not work. 

 
                    \begin{figure}[H]
                        \centering
                        \includegraphics[width=0.5\textwidth]{Solution/AIME_2024_I_11_Solution_Figure_6.png}
                    \end{figure}
                    

 If we have one adjacent pair and two separate ones that are $2$ away from each other, we can flip the diagram by $4$ vertices. There are $8$ of these cases. 

 
                    \begin{figure}[H]
                        \centering
                        \includegraphics[width=0.5\textwidth]{Solution/AIME_2024_I_11_Solution_Figure_7.png}
                    \end{figure}
                    

 Finally, if the red and blues alternate, we can simply shift the diagram by a single vertex to satisfy the question. Thus, all of these cases work, and we have $2$ subcases. 

 
                    \begin{figure}[H]
                        \centering
                        \includegraphics[width=0.5\textwidth]{Solution/AIME_2024_I_11_Solution_Figure_8.png}
                    \end{figure}
                    

 There can not be more than $4$ blues, so we are done. 

 Our total is $1+8+28+56+8+4+8+2=115$. There are $2^8=256$ possible colorings, so we have $\dfrac{115}{256}$ and our answer is $115+256=\boxed{371}$.

 





\textbf{Solution 2}

Let $r$ be the number of red vertices and $b$ be the number of blue vertices, where $r+b=8$. By the Pigeonhole Principle, $r\geq{b} \Longrightarrow b\leq4$ if a configuration is valid.

 We claim that if $b\leq3$, then any configuration is valid. We attempt to prove by the following:

 If there are 
 \[b\in{0,1,2}\] vertices, then intuitively any configuration is valid. For $b=3$, we do cases:

 If all the vertices in $b$ are non-adjacent, then simply rotating once in any direction suffices. If there are $2$ adjacent vertices, then WLOG let us create a set $\{b_1,b_2,r_1\cdots\}$ where the third $b_3$ is somewhere later in the set. If we assign the set as $\{1,2,3,4,5,6,7,8\}$ and $b_3\leq4$, then intuitively, rotating it $4$ will suffice. If $b_3=5$, then rotating it by 2 will suffice. Consider any other $b_3>5$ as simply a mirror to a configuration of the cases.

 Therefore, if $b\leq3$, then there are $\sum_{i=0}^{3}{\binom{8}{i}}=93$ ways. We do count the [i]degenerate[/i] case.

 Now if $b=4$, we do casework on the number of adjacent vertices.0 adjacent: $\{b_1,r_1,b_2,r_2\cdots{r_4}\}$. There are 4 axes of symmetry so there are only $\frac{8}{4}=2$ rotations of this configuration.

 1 adjacent: WLOG $\{b_1,b_2\cdots{b_3}\cdots{b_4}\}$ where $b_4\neq{8}$. Listing out the cases and trying, we get that $b_3=4$ and $b_4=7$ is the only configuration. There are $8$ ways to choose $b_1$ and $b_2$ and the rest is set, so there are $8$ ways.

 2 adjacent: We can have WLOG $\{b_1,b_2\cdots{b_3},b_4\}$ or $\{b_1,b_2,b_3\cdots\}$ where $b_4\neq{8}$. The former yields the case $b_3=5$ and $b_4=6$ by simply rotating it 2 times. The latter yields none. There are 2 axes of symmetry so there are $\frac{8}{2}=4$ configurations.

 3 adjacent: WLOG $\{b_1,b_2,b_3,b_4\cdots\}$ which intuitively works. There are $8$ configurations here as $b_1$ can is unique.

 In total, $b=4$ yields $2+8+4+8=22$ configurations.

 There are $22+93=115$ configurations in total. There are $2^8=256$ total cases, so the probability is $\frac{115}{256}$. Adding them up, we get $115+256=\boxed{371}$.

 




	\item (\textit{AIME I 2024}) Define $f(x)=|| x|-\frac{1}{2}|$ and $g(x)=|| x|-\frac{1}{4}|$. Find the number of intersections of the graphs of $y=4g(f(\sin (2 \pi x)))$ and $x=4g(f(\cos (3 \pi y)))$.


 
 \textbf{Answer: }385



\textbf{Solution 1}

If we graph $4g(f(x))$, we see it forms a sawtooth graph that oscillates between $0$ and $1$ (for values of $x$ between $-1$ and $1$, which is true because the arguments are between $-1$ and $1$). Thus by precariously drawing the graph of the two functions in the square bounded by $(0,0)$, $(0,1)$, $(1,1)$, and $(1,0)$, and hand-counting each of the intersections, we get $\boxed{385}$

 \textbf{Remark} While this solution might seem unreliable (it probably is), the only parts where counting the intersection might be tricky is near $(1,1)$. Make sure to count them as two points and not one, or you'll get $384$.

 \textbf{Note 1} The answer should be 385 since there are 16 intersections in each of 24 smaller boxes of dimensions 1/6 x 1/4 and then another one at the corner (1,1).

 





\textbf{Solution 2}

We will denote $h(x)=4g(f(x))$ for simplicity. Denote $p(x)$ as the first equation and $q(y)$ as the graph of the second. We notice that both $f(x)$ and $g(x)$ oscillate between 0 and 1. The intersections are thus all in the square $(0,0)$, $(0,1)$, $(1,1)$, and $(1,0)$. Every $p(x)$ wave going up and down crosses every $q(y)$ wave. Now, we need to find the number of times each wave touches 0 and 1.

 
We notice that $h(x)=0$ occurs at $x=-\frac{3}{4}, -\frac{1}{4}, \frac{1}{4}, \frac{3}{4}$, and $h(x)=1$ occurs at $x=-1, -\frac{1}{2}, 0,\frac{1}{2},1$. A sinusoid passes through each point twice during each period, but it only passes through the extrema once. $p(x)$ has 1 period between 0 and 1, giving 8 solutions for $p(x)=0$ and 9 solutions for $p(x)=1$, or 16 up and down waves. $q(y)$ has 1.5 periods, giving 12 solutions for $q(y)=0$ and 13 solutions for $q(y)=1$, or 24 up and down waves. This amounts to $16\cdot24=384$ intersections.

 
However, we have to be very careful when counting around $(1, 1)$. At this point, $q(y)$ has an infinite downwards slope and $p(x)$ is slanted, giving us an extra intersection; thus, we need to add 1 to our answer to get $\boxed{385}$.

 





\textbf{Solution 3 (Rigorous analysis of why there are two solutions near (1,1))}

We can easily see that only $x, y \in \left[0,1 \right]$ may satisfy both functions.We call function $y = 4g \left( f \left( \sin \left( 2 \pi x \right) \right) \right)$ as Function 1 and function $x = 4g \left( f \left( \cos \left( 3 \pi y \right) \right) \right)$ as Function 2.

 For Function 1, in each interval $\left[ \frac{i}{4} , \frac{i+1}{4} \right]$ with $i \in \left\{ 0, 1, \cdots, 3 \right\}$, Function 1's value oscillates between 0 and 1. It attains 1 at $x = \frac{i}{4}$, $\frac{i+1}{4}$ and another point between these two.Between two consecutive points whose functional values are 1, the function first decreases from 1 to 0 and then increases from 0 to 1. So the graph of this function in this interval consists of 4 monotonic pieces.

 For Function 2, in each interval $\left[ \frac{i}{6} , \frac{i+1}{6} \right]$ with $i \in \left\{ 0, 1, \cdots, 5 \right\}$, Function 2's value oscillates between 0 and 1. It attains 1 at $y = \frac{i}{6}$, $\frac{i+1}{6}$ and another point between these two.Between two consecutive points whose functional values are 1, the function first decreases from 1 to 0 and then increases from 0 to 1.So the graph of this function in this interval consists of 4 monotonic curves.

 Consider any region $\left[ \frac{i}{4} , \frac{i+1}{4} \right] \times \left[ \frac{j}{6} , \frac{j+1}{6} \right]$ with $i \in \left\{ 0, 1, \cdots, 3 \right\}$ and $j \in \left\{0, 1, \cdots , 5 \right\}$ but $\left( i, j \right) \neq \left( 3, 5 \right)$.Both functions have four monotonic pieces.Because Function 1's each monotonic piece can take any value in  $\left[ \frac{j}{6} , \frac{j+1}{6} \right]$ and Function 2' each monotonic piece can take any value in  $\left[ \frac{i}{4} , \frac{i+1}{4} \right]$, Function 1's each monotonic piece intersects with Function 2's each monotonic piece. Therefore, in the interval $\left[ \frac{i}{4} , \frac{i+1}{4} \right] \times \left[ \frac{j}{6} , \frac{j+1}{6} \right]$, the number of intersecting points is $4 \cdot 4 = 16$.

 Next, we prove that if an intersecting point is on a line $x = \frac{i}{4}$ for $i \in \left\{ 0, 1, \cdots, 4 \right\}$, then this point must be $\left( 1, 1 \right)$.

 For $x = \frac{i}{4}$, Function 1 attains value 1. For Function 2, if $y = 1$, then $x = 1$. Therefore, the intersecting point is $\left( 1, 1 \right)$.

 Similarly, we can prove that if an intersecting point is on a line $y = \frac{i}{6}$ for $i \in \left\{ 0, 1, \cdots, 6 \right\}$, then this point must be $\left( 1, 1 \right)$.

 Therefore, in each region $\left[ \frac{i}{4} , \frac{i+1}{4} \right] \times \left[ \frac{j}{6} , \frac{j+1}{6} \right]$ with $i \in \left\{ 0, 1, \cdots, 3 \right\}$ and $j \in \left\{0, 1, \cdots , 5 \right\}$ but $\left( i, j \right) \neq \left( 3, 5 \right)$, all 16 intersecting points are interior.That is, no two regions share any common intersecting point.

 Next, we study region  $\left[ \frac{3}{4} , 1 \right] \times \left[ \frac{5}{6} , 1 \right]$.Consider any pair of monotonic pieces, where one is from Function 1 and one is from Function 2, except the pair of two monotonic pieces from two functions that attain $\left( 1 , 1 \right)$.Two pieces in each pair intersects at an interior point on the region.So the number of intersecting points is $4 \cdot 4 - 1 = 15$.

 Finally, we compute the number intersection points of two functions' monotonic pieces that both attain $\left( 1, 1 \right)$.

 One trivial intersection point is $\left( 1, 1 \right)$.Now, we study whether they intersect at another point.

 Define $x = 1 - x'$ and $y = 1 - y'$.Thus, for positive and sufficiently small $x'$ and $y'$, Function 1 is reduced to
 \[ y' = 4 \sin 2 \pi x' \hspace{1cm} (1) \]and Function 2 is reduced to
 \[ x' = 4 \left( 1 - \cos 3 \pi y' \right) . \hspace{1cm} (2) \]

 Now, we study whether there is a non-zero solution.Because we consider sufficiently small $x'$ and $y'$, to get an intuition and quick estimate, we do approximations of the above equations.

 Equation (1) is approximated as
 \[ y' = 4 \cdot 2 \pi x' \]and Equation (2) is approximated as
 \[ x' = 2 \left( 3 \pi y' \right)^2  \]

 To solve these equations, we get $x' = \frac{1}{8^2 \cdot 18 \pi^4}$ and $y' = \frac{1}{8 \cdot 18 \pi^3}$.Therefore, two functions' two monotonic pieces that attain $\left( 1, 1 \right)$ have two intersecting points.

 Putting all analysis above, the total number of intersecting points is $16 \cdot 4 \cdot 6 + 1 = \boxed{\textbf{(385) }}$.

 


 \textbf{Drawing of two solutions near $\left( 1, 1 \right)$} \href{https://artofproblemsolving.com/wiki/index.php/File:2024\_AIME\_I\_Problem\_12,\_two\_solutions\_near\_(1,1).png}{https://artofproblemsolving.com/wiki/index.php/File:2024\_AIME\_I\_Problem\_12,\_two\_solutions\_near\_(1,1).png}

 


 




	\item (\textit{AIME I 2024}) Let $p$ be the least prime number for which there exists an integer $n$ such that $n^{4}+1$ is divisible by $p^{2}$. Find the least positive integer $m$ such that $m^{4}+1$ is divisible by $p^{2}$.


 
 \textbf{Answer: }110



\textbf{Solution 1}

If \(p=2\), then \(4\mid n^4+1\) for some integer \(n\). But \(\left(n^2\right)^2\equiv0\) or \(1\pmod4\), so it is impossible. Thus \(p\) is an odd prime.

 There are two ways to do the first part:

 \textbf{First Method using FLT.} For integer \(n\) such that \(p^2\mid n^4+1\), we have \(p\mid n^4+1\), hence \(p\nmid n^4-1\), but \(p\mid n^8-1\). By \href{https://artofproblemsolving.com/wiki/index.php?title=Fermat\%27s\_Little\_Theorem}{Fermat's Little Theorem}, \(p\mid n^{p-1}-1\), so

 \begin{equation*}p\mid\gcd\left(n^{p-1}-1,n^8-1\right)=n^{\gcd(p-1,8)}-1.\end{equation*}

 Here, \(\gcd(p-1,8)\) mustn't be divide into \(4\) or otherwise \(p\mid n^{\gcd(p-1,8)}-1\mid n^4-1\), which contradicts. So \(\gcd(p-1,8)=8\), and so \(8\mid p-1\). The smallest such prime is clearly \(p=17=2\times8+1\).

 
\textbf{Second Method using orders.} Notice that the order modulo $p^2$ of $n^4$ must be $8$, so $8\mid\varphi(p^2)=p(p-1)$. Hence $8\mid p-1$, so $p\equiv 1\pmod{8}$. The smallest such prime is $p=17$. ~eevee9406

 
So we have to find the smallest positive integer \(m\) such that \(17\mid m^4+1\). We first find the remainder of \(m\) divided by \(17\) by doing

 \begin{array}{|c|cccccccccccccccc|}\hline\vphantom{\frac11}x\bmod{17}&1&2&3&4&5&6&7&8&9&10&11&12&13&14&15&16\\\hline\vphantom{\dfrac11}\left(x^4\right)+1\bmod{17}&2&0&14&2&14&5&5&0&0&5&5&14&2&14&0&2\\\hline\end{array}

 So \(m\equiv\pm2\), \(\pm8\pmod{17}\). If \(m\equiv2\pmod{17}\), let \(m=17k+2\), by the binomial theorem,

 \begin{align*}0&\equiv(17k+2)^4+1\equiv\mathrm {4\choose 1}(17k)(2)^3+2^4+1=17(1+32k)\pmod{17^2}\\[3pt]\implies0&\equiv1+32k\equiv1-2k\pmod{17}.\end{align*}


 So the smallest possible \(k=9\), and \(m=155\).

 If \(m\equiv-2\pmod{17}\), let \(m=17k-2\), by the binomial theorem,

 \begin{align*}0&\equiv(17k-2)^4+1\equiv\mathrm {4\choose 1}(17k)(-2)^3+2^4+1=17(1-32k)\pmod{17^2}\\[3pt]\implies0&\equiv1-32k\equiv1+2k\pmod{17}.\end{align*}


 So the smallest possible \(k=8\), and \(m=134\).

 If \(m\equiv8\pmod{17}\), let \(m=17k+8\), by the binomial theorem,

 \begin{align*}0&\equiv(17k+8)^4+1\equiv\mathrm {4\choose 1}(17k)(8)^3+8^4+1=17(241+2048k)\pmod{17^2}\\[3pt]\implies0&\equiv241+2048k\equiv3+8k\pmod{17}.\end{align*}


 So the smallest possible \(k=6\), and \(m=110\).

 If \(m\equiv-8\pmod{17}\), let \(m=17k-8\), by the binomial theorem,

 \begin{align*}0&\equiv(17k-8)^4+1\equiv\mathrm {4\choose 1}(17k)(-8)^3+8^4+1=17(241-2048k)\pmod{17^2}\\[3pt]\implies0&\equiv241-2048k\equiv3+9k\pmod{17}.\end{align*}


 So the smallest possible \(k=11\), and \(m=179\).

 In conclusion, the smallest possible \(m\) is \(\boxed{110}\).

 <font size="2">Solution by Quantum-Phantom</font>

 





\textbf{Solution 2 (Hensel's Lemma)}

We know, from Solution 1, that by Fermat's Little Theorem, $p\equiv 1\pmod{8}$, and that the smallest $p$ satisfying such is $17$. Thus, we have $p^2=289$ and we are attempting to lift modulo $17$ to modulo $289$.

 Given in the problem, it is established that $p^2\mid m^4 +1$ holds true, so we can guarantee that $m^4\equiv 1\pmod{289}$. Testing for roots $r_0$ modulo $17$ gives us the solutions $r_0= \pm 2$, $\pm 8$, which satisfy $r_0^4\equiv -1 \pmod{17}$. 

 Note that we can now use Hensel's Lemma, by defining a function $f(x)=x^4+1$, giving $f'(x)=4x^3$. We know firstly that, though $f(r_0)$ for all roots $r_0$ we found is congruent to $0\pmod{17}$, this does not hold true for $f'(r_0)$, so we guarantee that there is one lift by Hensel's Lemma to $289$ that we can perform, defined for root $r_1$ that $r_1 = r_0-\frac{f(r_0)}{f'(r_0)}\pmod{p^k}$. We now break our work into four cases:

 
            \begin{enumerate}
                \item $r_0 = 2$. $f(2)=17$ and $f'(2)=32$, so $r_1=2-\frac{17}{32}\pmod{289}$, and we need to find the inverse of $32\equiv 15$ modulo $17$. By the extended Euclidean algorithm, we find that $8\cdot 15\equiv 1 \pmod{17}$, and thus $r_1=2-17\cdot 8\equiv 155 \pmod{289}$.\item $r_0 = -2$, or equivalent, $15$. $f(-2)=17$ and $f'(-2)=-32$, so $r_1=-2+\frac{17}{32}\pmod{289}$. We already found the inverse of $32$ modulo $17$, so this case has $r_1=-1+17\cdot 8\equiv 134 \pmod{289}$.\item $r_0 = 8$. $f(8)=4097$ and $f'(8)=2048$, so $r_1=-2+\frac{4097}{2048}\pmod{289}$, and we need to find the inverse of $32\equiv 8$ modulo $17$. We already know the inverse of $8$ is $15$, so we end up with $r_1=8-4097\cdot 15\equiv 110\pmod{289}$.\item $r_0 = -8$, or equivalent, $9$. $f(-8)=4097$ and $f'(-8)=-2048$, so $r_1=-8+\frac{4097}{2048}\pmod{289}$. We already found the inverse of $2048$ modulo $17$, so this case has $r_1=-8+4097\cdot 15\equiv 179\pmod{289}$.
            \end{enumerate}
            Thus, out of our four values for $m$, the smallest is $\boxed{110}$. Solution by \href{https://artofproblemsolving.com/wiki/index.php?title=User:Juwushu}{juwushu}.

 





\textbf{Solution 3 (Easy, given specialized knowledge)}

Note that $n^4 + 1 \equiv 0 \pmod{p}$ means $\text{ord}_{p}(n) = 8 \mid p-1.$ The smallest prime that does this is $17$ and $2^4 + 1 = 17$ for example. Now let $g$ be a primitive root of $17^2.$ The satisfying $n$ are of the form, $g^{\frac{p(p-1)}{8}}, g^{3\frac{p(p-1)}{8}}, g^{5\frac{p(p-1)}{8}}, g^{7\frac{p(p-1)}{8}}.$ So if we find one such $n$, then all $n$ are $n, n^3, n^5, n^7.$ Consider the $2$ from before. Note $17^2 \mid 2^{4 \cdot 17} + 1$ by LTE. Hence the possible $n$ are, $2^{17}, 2^{51}, 2^{85}, 2^{119}.$ Some modular arithmetic yields that $2^{51} \equiv \boxed{110}$ is the least value. 

 





\textbf{Solution 4 (Unfinished)}

We work in the ring \(\mathbb Z/289\mathbb Z\) and use the formula
 \[\sqrt[4]{-1}=\pm\sqrt{\frac12}\pm\sqrt{-\frac12}.\]Since \(-\frac12=144\), the expression becomes \(\pm12\pm12i\), and it is easily calculated via Hensel that \(i=38\), thus giving an answer of \(\boxed{110}\).

 





\textbf{Solution 5}

Alternative Beginning to Solution 1.

 Note $n^4 + 1 \equiv 0 \pmod{p^2}$ implies $n^4 + 1 \equiv 0 \pmod{p}$, so $n^4 \equiv p-1 \pmod{p}$. Also, $p$ must satisfy $p^2 - 1 \pmod{p} \equiv n^4$. The first constraint implies the smallest values of $p$ are $p = 2,17$. However, if $p = 2$, then $p^2 - 1 \equiv 3 \pmod{4}$. However, all quartic residues (as in solution 1), are $0, 1 \pmod{4}$, so 2 is impossible. Try $p = 17$. We see that $n = 2$ satisfies the two equations, as $16 \equiv 17-1 \pmod{17}$, and $289 - 1 \pmod{289} \equiv n^4$ is a possible quartic residue, as $288 \pmod{289} \equiv -1 \pmod{289}$ which has a solution modulo 17, specifically when $n = 2$ in which we let $x = 2+17k$, which when raised to the fourth taken modulo 289, we obtain $x^4 \equiv 16 + 544k \pmod{289}$. Reduce 544 and obtain $x^4 \equiv 16 + 255k \pmod{289}$. Now, assume this takes the form $16 + 255k \equiv -1 \equiv 288 \pmod{289}$. Then, $255k \equiv 272 \pmod{289}$ must have solutions. Divide by 17, and obtain $15k \equiv 16 \pmod{17}$. The inverse of $15k$ is 8, so $k \equiv 9 \pmod{17}$. Because solutions for $k$ exist, that directly implies 288 is a valid quartic residue, so $p = 17$ is our smallest prime.

 Continue as in solution 1 to obtain $\boxed{110}$.

 


 




	\item (\textit{AIME I 2024}) Let $ABCD$ be a tetrahedron such that $AB = CD = \sqrt{41}$, $AC = BD = \sqrt{80}$, and $BC = AD = \sqrt{89}$. There exists a point $I$ inside the tetrahedron such that the distances from $I$ to each of the faces of the tetrahedron are all equal. This distance can be written in the form $\frac{m \sqrt{n}}{p}$, when $m$, $n$, and $p$ are positive integers, $m$ and $p$ are relatively prime, and $n$ is not divisible by the square of any prime. Find $m+n+p$.


 
 \textbf{Answer: }104



\textbf{Solution 1}

Notice that \(41=4^2+5^2\), \(89=5^2+8^2\), and \(80=8^2+4^2\), let \(A~(0,0,0)\), \(B~(4,5,0)\), \(C~(0,5,8)\), and \(D~(4,0,8)\). Then the plane \(BCD\) has a normal

 \begin{equation*}\mathbf n:=\frac14\overrightarrow{BC}\times\overrightarrow{CD}=\frac14\begin{pmatrix}-4\\0\\8\end{pmatrix}\times\begin{pmatrix}4\\-5\\0\end{pmatrix}=\begin{pmatrix}10\\8\\5\end{pmatrix}.\end{equation*}

 Hence, the distance from \(A\) to plane \(BCD\), or the height of the tetrahedron, is

 \begin{equation*}h:=\frac{\mathbf n\cdot\overrightarrow{AB}}{|\mathbf n|}=\frac{10\times4+8\times5+5\times0}{\sqrt{10^2+8^2+5^2}}=\frac{80\sqrt{21}}{63}.\end{equation*}

 Each side of the tetrahedron has the same area due to congruency by "S-S-S", and we call it \(S\). Then by the volume formula for pyramids,

 \begin{align*}\frac13Sh&=V_{D\text-ABC}=V_{I\text-ABC}+V_{I\text-BCD}+V_{I\text-CDA}+V_{I\text-DAB}\\&=\frac13Sr\cdot4.\end{align*}


 Hence, \(r=\frac h4=\frac{20\sqrt{21}}{63}\), and so the answer is \(20+21+63=\boxed{104}\).

 <font size="2">Solution by Quantum-Phantom</font>

 





\textbf{Solution 2}


                    \begin{figure}[H]
                        \centering
                        \includegraphics[width=0.5\textwidth]{Solution/AIME_2024_I_14_Solution_Figure_1.png}
                    \end{figure}
                    

 
Inscribe tetrahedron $ABCD$ in an rectangular prism as shown above.

 By the Pythagorean theorem, we note

 
 \[OA^2 + OB^2 = AB^2 = 41,\]
 \[OA^2 + OC^2 = AC^2 = 80, \text{and}\]
 \[OB^2 + OC^2 = BC^2 = 89.\]

 Solving yields $OA = 4, OB = 5,$ and $OC = 8.$

 Since each face of the tetrahedron is congruent, we know the point we seek is the center of the circumsphere of $ABCD.$ We know all rectangular prisms can be inscribed in a circumsphere, therefore the circumsphere of the rectangular prism is also the circumsphere of $ABCD.$

 We know that the distance from all $4$ faces must be the same, so we only need to find the distance from the center to plane $ABC$.

 Let $O = (0,0,0), A = (4,0,0), B = (0,5,0),$ and $C = (0,0,8).$ We obtain that the plane of $ABC$ can be marked as $\frac{x}{4} + \frac{y}{5} + \frac{z}{8} = 1,$ or $10x + 8y + 5z - 40 = 0,$ and the center of the prism is $(2,\frac{5}{2},4).$

 Using the Point-to-Plane distance formula, our distance is

 
 \[d = \frac{|10\cdot 2 + 8\cdot \frac{5}{2} + 5\cdot 4 - 40|}{\sqrt{10^2 + 8^2 + 5^2}} = \frac{20}{\sqrt{189}} = \frac{20\sqrt{21}}{63}.\]

 Our answer is $20 + 21 + 63 = \boxed{104}.$

 - \href{https://artofproblemsolving.com/wiki/index.php/User:Spectraldragon8}{spectraldragon8}

 





\textbf{Solution 3(Formula Abuse)}

We use the formula for the volume of iscoceles tetrahedron. $V = \sqrt{(a^2 + b^2 - c^2)(b^2 + c^2 - a^2)(a^2 + c^2 - b^2)/72}$

 Note that all faces have equal area due to equal side lengths. By Law of Cosines, we find 
 \[\cos{\angle ACB} = \frac{80 + 89 - 41}{2\sqrt{80\cdot 89}}= \frac{16}{\sqrt{445}}.\]. 

 From this, we find 
 \[\sin{\angle ACB} = \sqrt{1-\cos^2{\angle ACB}} = \sqrt{1 - \frac{256}{445}} = \sqrt{\frac{189}{445}}\] and can find the area of $\triangle ABC$ as 
 \[A = \frac{1}{2} \sqrt{89\cdot 80}\cdot \sin{\angle ACB} = 6\sqrt{21}.\]

 Let $R$ be the distance we want to find. By taking the sum of (equal) volumes 
 \[[ABCI] + [ABDI] + [ACDI] + [BCDI] = V,\]  We have 
 \[V = \frac{4AR}{3}.\] Plugging in and simplifying, we get $R = \frac{20\sqrt{21}}{63}$ for an answer of $20 + 21 + 63 = \boxed{104}$

 





\textbf{Solution 4}

Let $AH$ be perpendicular to $BCD$ that meets this plane at point $H$.Let $AP$, $AQ$, and $AR$ be heights to lines $BC$, $CD$, and $BD$ with feet $P$, $Q$, and $R$, respectively.

 We notice that all faces are congruent. Following from Heron's formula, the area of each face, denoted as $A$, is $A = 6 \sqrt{21}$.

 Hence, by using this area, we can compute $AP$, $AQ$ and $AR$.We have $AP = \frac{2 A}{BC} = \frac{2A}{\sqrt{89}}$, $AQ = \frac{2 A}{CD} = \frac{2A}{\sqrt{41}}$, and $AR = \frac{2 A}{BC} = \frac{2A}{\sqrt{80}}$.

 Because $AH \perp BCD$, we have $AH \perp BC$. Recall that $AP \perp BC$.Hence, $BC \perp APH$. Hence, $BC \perp HP$.

 Analogously, $CD \perp HQ$ and $BD \perp HR$.

 We introduce a function $\epsilon \left( l \right)$ for $\triangle BCD$ that is equal to 1 (resp. -1) if point $H$ and the opposite vertex of side $l$ are on the same side (resp. opposite sides) of side $l$.

 The area of $\triangle BCD$ is

 \begin{align*}A & = \epsilon_{BC} {\textrm Area} \ \triangle HBC+ \epsilon_{CD} {\textrm Area} \ \triangle HCD+ \epsilon_{BD} {\textrm Area} \ \triangle HBD \\& = \frac{1}{2} \epsilon_{BC} BC \cdot HP+ \frac{1}{2} \epsilon_{CD} CD \cdot HQ +\frac{1}{2} \epsilon_{BD} BD \cdot HR \\& = \frac{1}{2} \epsilon_{BC} BC \cdot \sqrt{AP^2 - AH^2}+ \frac{1}{2} \epsilon_{CD} CD \cdot \sqrt{AQ^2 - AH^2} \\& \quad + \frac{1}{2} \epsilon_{BD} CD \cdot \sqrt{AR^2 - AH^2} . \hspace{1cm} (1)\end{align*}


 Denote $B = 2A$.The above equation can be organized as

 \begin{align*}B & = \epsilon_{BC} \sqrt{B^2 - 89 AH^2}+ \epsilon_{CD} \sqrt{B^2 - 41 AH^2} \\& \quad + \epsilon_{BD} \sqrt{B^2 - 80 AH^2} .\end{align*}


 This can be further reorganized as

 \begin{align*}B - \epsilon_{BC} \sqrt{B^2 - 89 AH^2}& = \epsilon_{CD} \sqrt{B^2 - 41 AH^2}+ \epsilon_{BD} \sqrt{B^2 - 80 AH^2} .\end{align*}


 Taking squares on both sides and reorganizing terms, we get

 \begin{align*}& 16 AH^2 - \epsilon_{BC} B \sqrt{B^2 - 89 AH^2} \\& = \epsilon_{CD} \epsilon_{BD}\sqrt{\left( B^2 - 41 AH^2 \right) \left( B^2 - 80 AH^2 \right)} .\end{align*}


 Taking squares on both sides and reorganizing terms, we get

 
 \[ - \epsilon_{BC} 2 B \sqrt{B^2 - 89 AH^2} = - 2 B^2 + 189 AH^2 . \]

 Taking squares on both sides, we finally get

 \begin{align*}AH & = \frac{20B}{189} \\& = \frac{40A}{189}.\end{align*}


 Now, we plug this solution to Equation (1). We can see that $\epsilon_{BC} = -1$, $\epsilon_{CD} = \epsilon_{BD} = 1$.This indicates that $H$ is out of $\triangle BCD$. To be specific, $H$ and $D$ are on opposite sides of $BC$, $H$ and $C$ are on the same side of $BD$, and $H$ and $B$ are on the same side of $CD$. 

 Now, we compute the volume of the tetrahedron $ABCD$, denoted as $V$. We have $V = \frac{1}{3} A \cdot AH = \frac{40 A^2}{3 \cdot 189}$.

 Denote by $r$ the inradius of the inscribed sphere in $ABCD$.Denote by $I$ the incenter.Thus, the volume of $ABCD$ can be alternatively calculated as

 \begin{align*}V & = {\textrm Vol} \ IABC + {\textrm Vol} \ IACD + {\textrm Vol} \ IABD + {\textrm Vol} \ IBCD \\& = \frac{1}{3} r \cdot 4A .\end{align*}


 From our two methods to compute the volume of $ABCD$ and equating them, we get

 \begin{align*}r & = \frac{10A}{189} \\& = \frac{20 \sqrt{21}}{63} .\end{align*}


 Therefore, the answer is $20 + 21 + 63 = \boxed{\textbf{(104) }}$.

 





\textbf{Solution 5(A quicker method to compute the height from  to plane )}

We put the solid to a 3-d coordinate system. Let $B = \left( 0, 0, 0 \right)$, $D = \left( \sqrt{80}, 0, 0 \right)$.We put $C$ on the $x-o-y$ plane.Now, we compute the coordinates of $C$.

 Applying the law of cosines on $\triangle BCD$, we get $\cos \angle CBD = \frac{4}{\sqrt{41 \cdot 5}}$.Thus, $\sin \angle CBD = \frac{3 \sqrt{21}}{\sqrt{41 \cdot 5}}$.Thus, $C = \left( \frac{4}{\sqrt{5}} , \frac{3 \sqrt{21}}{\sqrt{5}} , 0 \right)$.

 Denote $A = \left( x, y , z \right)$ with $z > 0$.

 Because $AB = \sqrt{89}$, we have$x^2 + y^2 + z^2 = 89$

 Because $AD = \sqrt{41}$, we have

 
 \[\left( x - \sqrt{80} \right)^2 + y^2 + z^2 = 41 \hspace{1cm} (2)\]

 Because $AC = \sqrt{80}$, we have

 
 \[\left( x - \frac{4}{\sqrt{5}} \right)^2 + \left( y - \frac{3 \sqrt{21}}{\sqrt{5}} \right)^2  + z^2 = 80 \hspace{1cm} (3)\]

 Now, we compute $x$, $y$ and $z$.

 Taking $(1)-(2)$, we get
 \[2 \sqrt{80} x = 128 .\]

 Thus, $x = \frac{16}{\sqrt{5}}$.

 Taking $(1) - (3)$, we get
 \[2 \cdot \frac{4}{\sqrt{5}} x + 2 \cdot \frac{3 \sqrt{21}}{\sqrt{5}} y = 50 .\]

 Thus, $y = \frac{61}{3 \sqrt{5 \cdot 21}}$.

 Plugging $x$ and $y$ into Equation (1), we get $z = \frac{80 \sqrt{21}}{63}$.

 





\textbf{Solution 6 (Different Perspective)}

Consider the following construction of the tetrahedron. Place $AB$ on the floor. Construct an isosceles vertical triangle with $AB$ as its base and $M$ as the top vertex. Place $CD$ on the top vertex parallel to the ground with midpoint $M.$ Observe that $CD$ can rotate about its midpoint. At a certain angle, we observe that the lengths satisfy those given in the problem. If we project $AB$ onto the plane of $CD$, let the minor angle $\theta$ be this discrepancy.

 By Median formula or Stewart's theorem, $AM = \frac{1}{2}\sqrt{2AC^2 + 2AD^2 - CD^2} = \frac{3\sqrt{33}}{2}.$ Consequently the area of $\triangle AMB$ is $\frac{\sqrt{41}}{2} \left (\sqrt{(\frac{3\sqrt{33}}{2})^2 - (\frac{\sqrt{41}}{2})^2} \right ) = 4\sqrt{41}.$ Note the altitude $8$ is also the distance between the parallel planes containing $AB$ and $CD.$

 By Distance Formula,

 
 \begin{align*} (\frac{\sqrt{41}}{2} - \frac{1}{2}CD \cos{\theta})^2 + (\frac{1}{2}CD \sin{\theta}^2) + (8)^2 &= AC^2 = 80 \\ (\frac{\sqrt{41}}{2} + \frac{1}{2}CD \cos{\theta})^2 + (\frac{1}{2}CD \sin{\theta}^2) + (8)^2 &= AD^2 = 89 \\ \implies CD \cos{\theta} \sqrt{41} &= 9 \\ \sin{\theta} &= \sqrt{1 - (\frac{9}{41})^2} = \frac{40}{41}. \end{align*}


 Then the volume of the tetrahedron is given by $\frac{1}{3} [AMB] \cdot CD \sin{\theta} = \frac{160}{3}.$

 The volume of the tetrahedron can also be segmented into four smaller tetrahedrons from $I$ w.r.t each of the faces. If $r$ is the inradius, i.e the distance to the faces, then $\frac{1}{3} r([ABC] + [ABD] + [ACD] + [BCD])$ must the volume. Each face has the same area by SSS congruence, and by Heron's it is $\frac{1}{4}\sqrt{(a + b + c)(a + b - c)(c + (a-b))(c -(a - b))} = 6\sqrt{21}.$

 Therefore the answer is, $\dfrac{3 \frac{160}{3}}{24 \sqrt{21}} = \frac{20\sqrt{21}}{63} \implies \boxed{104}.$

 





\textbf{Solution 7 (simplest way to calculate [ABCD])}

After finding that $ABCD$ can be inscribed in a $4\times5\times8$ box, note that $[ABCD]=\mathbf{box}-(\frac{1}{2})(\frac{4}{3})(4)(5)(8)=\frac{160}{3}$, where the third term describes the area of 4 congruent right tetrahedrons whose right angle is at the apex and whose apex lengths are 4, 5, and 8. Then proceed as in Solution 3.

 





\textbf{Solution 8 (fast)}

Observe that $ABCD$ can be inscribed in a $4\times5\times8$ box, which gives the coordinates of the points: we have $A=(8,0,0)$, $B=(8,5,4)$, $C=(0,0,4)$ and $D=(0,5,0)$. Recall the formula $V=\frac{1}{3}RS$, where $V$ is the volume of the tetrahedron, $S$ is the surface area, and $R$ is the insphere radius. Using the fact that the faces of $ABCD$ have equal area and that $V=\frac{1}{3}[ACD] \times h$, where $h$ is the distance from plane $ACD$ to $B$, we can rearrange to get $R=\frac{h}{4}$. We compute that the the equation of plane $ACD$ is $5x+8y+10z-40=0$, and applying point-to-plane gives $h=\frac{\vert 5(8)+8(5)+10(4)-40 \vert}{\sqrt{5^2+8^2+10^2}}=\frac{80}{\sqrt{189}} \implies R=\frac{20\sqrt{21}}{63} \implies \boxed{104}$.

 





\textbf{Solution 9 (Analytical Geometry)}

Like in Solution 1, we have $A = (0,0,0)$, $B = (4,5,0)$, $C = (0,5,8)$, $D = (4,0,8)$. If we let the inradius be $r$, then the volume of the tetrahedron is equal to $\frac{1}{3}[ABC]r + \frac{1}{3}[BCD]r + \frac{1}{3}[ABD]r + \frac{1}{3}[ACD]r = V$ by the formula $V = \frac{1}{3}Ah$. Also note that $[ABC] = [BCD] = [ABD] = [ACD]$ because they all have the same side lengths. So, if we let $[ABC] = A$, then $\frac{4}{3}Ar = V$, so $r = \frac{3V}{4A}$. Now, all we need to do is to find V and A. We note that the area of the triangle $[ABC] = |\frac{\overrightarrow{AB}\times\overrightarrow{AC}}{2}|$. Also, if we let the normal vector from $\Delta ABC$ to be $\bold{n}$, the height $h$ from $D$ to $\Delta ABC$ is $\overrightarrow{AD} \cdot \hat{\bold{n}} = \frac{\overrightarrow{AD} \cdot \bold{n}}{\bold{|n|}}$. Therefore, the volume of the tetrahedron is equal to $\frac{1}{3}[ABC]h = \frac{1}{3}|\frac{\overrightarrow{AB}\times\overrightarrow{AC}}{2}|\frac{\overrightarrow{AD} \cdot \bold{n}}{\bold{|n|}} = \frac{1}{6}|\overrightarrow{AB}\times\overrightarrow{AC}|\frac{\overrightarrow{AD} \cdot \bold{n}}{\bold{|n|}}$. Now, we notice that $\overrightarrow{AB}\times\overrightarrow{AC}$ is a normal vector of the $\Delta ABC$ because the cross product of two vectors is always perpendicular to both of the vectors, and both of the vectors lie in the plane of $\Delta ABC$. This means that we can let $\bold{n} = \overrightarrow{AB}\times\overrightarrow{AC}$. So, the volume of the tetrahedron is $V = \frac{1}{6}|\overrightarrow{AB}\times\overrightarrow{AC}|\frac{\overrightarrow{AD} \cdot (\bold{\overrightarrow {AB}\times\overrightarrow{AC}})}{\bold{|\overrightarrow{AB}\times\overrightarrow{AC}|}} = \frac{1}{6}\overrightarrow{AD} \cdot (\bold{\overrightarrow{AB}\times\overrightarrow{AC}})$. The expression $\overrightarrow{AD} \cdot (\bold{\overrightarrow{AB}\times\overrightarrow{AC}})$ is more commonly known as a box product and can be used to find the volume of any parallelepiped. We know that $\overrightarrow{AB} = (4,5,0)$, $\overrightarrow{AC} = (0,5,8)$, and $\overrightarrow{AD} = (4,0,8)$. Now, $\overrightarrow{AB} \times \overrightarrow{AC} = (8*5-5*0, 0*0-8*4, 4*5-5*0) = (40, -32, 20)$. So, $\overrightarrow{AD} \cdot (\bold{\overrightarrow{AB}\times\overrightarrow{AC}}) = (4,0,8) \cdot (40,-32,20) = 4*40 + 0*-32 + 8*20 = 320$. So, $V = \frac{1}{6}\overrightarrow{AD} \cdot (\bold{\overrightarrow{AB}\times\overrightarrow{AC}}) = \frac{320}{6} = \frac{160}{3}$. Now $[ABC]$ is just equal to $\frac{1}{2}|\overrightarrow{AB} \times \overrightarrow{AC}|$ as it is half the parallelogram. So, we have $A = \frac{1}{2}|\overrightarrow{AB} \times \overrightarrow{AC}| = \frac{1}{2}|(40,-32,20)| = \frac{1}{2}\sqrt{3024}$. So, $r = \frac{3V}{4A} = \frac{160}{2\sqrt{3024}}$. This expression can be simplified to $\frac{20\sqrt{21}}{63}$. So, our answer is $20 + 63 + 21 = \boxed{104}$.

 





\textbf{Solution 10 (simplest)}

Obviously, $ABCD$ is a \href{https://artofproblemsolving.com/wiki/index.php?title=Disphenoid}{disphenoid}, and it follows that its circumscribing rectangular prism with nonadjacent vertices $A$, $B$, $C$, and $D$ has a length, width, and height of $4$, $5$, and $8$. Thus, 
 \[[ABCD]=4\cdot 5\cdot 8\left(1-4\cdot\frac 16\right)=\frac{160}{3}.\] Since the distances from $I$ to the faces of $ABCD$ are all equal, let this common value be $x$ and WLOG it follows that

 \begin{align*}\frac 13x[ABC]&=[IABC],\\2x\sqrt{21}&=\frac 14\cdot\frac{160}{3},\\x&=\frac{20\sqrt{21}}{63}\\\end{align*}


 through Heron's formula, so the desired value is $\boxed{104}$.

 -pieMax2713

 


 




	\item (\textit{AIME I 2024}) Let $\mathcal{B}$ be the set of rectangular boxes with surface area $54$ and volume $23$. Let $r$ be the radius of the smallest sphere that can contain each of the rectangular boxes that are elements of $\mathcal{B}$. The value of $r^2$ can be written as $\frac{p}{q}$, where $p$ and $q$ are relatively prime positive integers. Find $p+q$.


 
 \textbf{Answer: }721



\textbf{Solution 1}

Observe that the "worst" possible box is one of the maximum possible length. By symmetry, the height and the width are the same in this antioptimal box. (If the height and width weren't the same, the extra difference between them could be used to make the length longer.) Thus, let the width and height be of length $a$ and the length be $L$. 

 We're given that the volume is $23$; thus, $a^2L=23$. We're also given that the surface area is $54=2\cdot27$; thus, $a^2+2aL=27$. 

 From the first equation, we can get $L=\dfrac{23}{a^2}$. We do a bunch of algebra: 

 \begin{align*}L&=\dfrac{23}{a^2} \\27&=a^2+2aL \\&=a^2+2a\left(\dfrac{23}{a^2}\right) \\&=a^2+\dfrac{46}a \\27a&=a^3+46 \\a^3-27a+46&=0. \\\end{align*}


 We can use the Rational Root Theorem and test a few values. It turns out that $a=2$ works. We use synthetic division to divide by $a-2$: 

 
                    \begin{figure}[H]
                        \centering
                        \includegraphics[width=0.5\textwidth]{Solution/AIME_2024_I_15_Solution_Figure_1.png}
                    \end{figure}
                    

 As we expect, the remainder is $0$, and we are left with the polynomial $x^2+2x-23$. We can now simply use the quadratic formula and find that the remaining roots are $\dfrac{-2\pm\sqrt{4-4(-23)}}2=\dfrac{-2\pm\sqrt{96}}2=\dfrac{-2\pm4\sqrt{6}}2=-1\pm2\sqrt6$. We want the smallest $a$ to maximize $L$, and it turns out that $a=2$ is in fact the smallest root. Thus, we let $a=2$. Substituting this into $L=\dfrac{23}{a^2}$, we find that $L=\dfrac{23}4$. However, this is not our answer! This is simply the length of the box; we want the radius of the sphere enclosing it. We know that the diameter of the sphere is the diagonal of the box, and the 3D Pythagorean Theorem can give us the space diagonal. Applying it, we find that the diagonal has length $\sqrt{2^2+2^2+\left(\dfrac{23}4\right)^2}=\sqrt{8+\dfrac{529}{16}}=\sqrt{\dfrac{128+529}{16}}=\dfrac{\sqrt{657}}4$. This is the diameter; we halve it to find the radius, $\dfrac{\sqrt{657}}8$. We then square this and end up with $\dfrac{657}{64}$, giving us an answer of $657+64=\boxed{721}$. 

 





\textbf{Solution 2 (constrained optimization with Lagrangian multiplier)}

Denote by $x$, $y$, $z$ the length, width, and height of a rectangular box.We have\begin{align*}xy + yz + zx & = 27  \hspace{1cm} (1) \\xyz & = 23 \hspace{1cm} (2)\end{align*}


 We have\begin{align*}4 r^2 & = x^2 + y^2 + z^2 \\& = \left( x + y + z \right)^2 - 2 \cdot \left( xy + yz + zx \right) \\& = \left( x + y + z \right)^2 - 54 .\end{align*}


 Therefore, we solve the following constrained optimization problem:\begin{align*}\max_{x,y,z} \ & x + y + z \\\mbox{subject to } & (1), (2) \end{align*}


 First, we prove that an optimal solution must have at least two out of $x$, $y$, $z$ that are the same.

 Denote by $\lambda$ and $\eta$ lagrangian multipliers of constraints (1) and (2), respectively. Consider the following Lagrangian:\begin{align*}\max_{x,y,z, \lambda, \eta} & x + y + z + \lambda \left( xy + yz + zx - 27 \right)+ \eta \left( xyz - 23 \right) .\end{align*}


 Taking first-order-condition with respect to $x$, $y$, $z$, respectively, we get\begin{align*}1 + \lambda \left( y + z \right) + \eta yz & = 0 \hspace{1cm} (3) \\1 + \lambda \left( z + x \right) + \eta zx & = 0 \hspace{1cm} (4) \\1 + \lambda \left( x + y \right) + \eta xy & = 0 \hspace{1cm} (5)\end{align*}


 Suppose there is an optimal solution with $x$, $y$, $z$ that are all distinct.

 Taking $(4)-(3)$, we get
 \[ \left( x - y \right) \left( \lambda + \eta z \right) = 0 . \]

 Because $x \neq y$, we have
 \[ \lambda + \eta z = 0 \hspace{1cm} (6) \]

 Analogously, we have\begin{align*}\lambda + \eta x & = 0 \hspace{1cm} (7) \end{align*}
 

 Taking $(6) - (7)$, we get $\eta \left( z - x \right) = 0$.Because $z \neq x$, we have $\eta = 0$. Plugging this into (6), we get $\lambda = 0$.

 However, the solution that $\lambda = \eta = 0$ is a contradiction with (3).Therefore, in an optimal solution, we cannot have $x$, $y$, and $z$ to be all distinct.

 W.L.O.G, in our remaining analysis, we assume an optimal solution satisfies $y = z$.

 Therefore, we need to solve the following two-variable optimization problem:\begin{align*}\max_{x,y} \ & x + 2y \\\mbox{subject to } & 2 xy + y^2 = 27 \\& xy^2 = 23\end{align*}


 Replacing $x$ with $y$ by using the constraint $xy^2 = 23$, we solve the following single-variable optimization problem:\begin{align*}\max_y \ &  \frac{23}{y^2} + 2y \hspace{1cm} (8) \\\mbox{subject to } & \frac{46}{y} + y^2 = 27 \hspace{1cm} (9)\end{align*}


 By solving (9), we get $y = 2$ and $-1 + 2 \sqrt{6}$.

 Plugging $y = 2$ into (8), we get $\frac{23}{y^2} + 2y = \frac{39}{4}$.

 Plugging $y = -1 + 2 \sqrt{6}$ into (8), we get $\frac{23}{y^2} + 2y = \frac{96 \sqrt{6} - 21}{23}$.

 We have $\frac{96 \sqrt{6} - 21}{23} < \frac{39}{4}$.Therefore, the maximum value of $x + y + z$ is $\frac{39}{4}$.

 Therefore,\begin{align*}r^2 & = \frac{1}{4} \left( \left( x + y + z \right)^2 - 54 \right) \\& = \frac{1}{4} \left( \left( \frac{39}{4} \right)^2 - 54 \right) \\& = \frac{657}{64} .\end{align*}


 Therefore, the answer is$657 + 64 = \boxed{\textbf{(721) }}$.

 





\textbf{Solution 3 (Vieta's Formula and Rational Root Theroem)}

First, let's list the conditions:Denote by $l$, $w$, $h$ the length, width, and height of a rectangular box.

 
 \[lwh=23\]\begin{align*}2(lw+wh+hl)&=54\\lw+wh+hl&=27.\end{align*}
Applying the Pythagorean theorem, we can establish that

 \begin{align*}(2r)^2&=(l^2+w^2+h^2)\\4r^2&=(l^2+w^2+h^2)\\4r^2&=(l+w+h)^2-2(lw+wh+hl)\\4r^2&=(l+w+h)^2-54.\end{align*}


 We can spot Vieta's formula hidden inside this equation and call this $m$. Now we have three equations:

 
 \[lwh=23\]
 \[lw+wh+hl)=27\] 
 \[l+w+h=m\]

 Let there be a cubic equation. $x^3+bx^2+cx+d=0$. Its roots are $l$, $w$ and $h$. We can use our formulas from before to derive $c$ and $d$.

 
 \[-b=l+w+h=m\]

 
 \[c=lw+wh+lh=27\]

 
 \[-d=lwh=23\]

 We can now rewrite the equation from before:

 $x^3-mx^2+27x-23=0$

 To find the maximum $r$ we need the maximum $m$. This only occurs when this equation has double roots illustrated with graph below.

 
                    \begin{figure}[H]
                        \centering
                        \includegraphics[width=0.5\textwidth]{Solution/AIME_2024_I_15_Solution_Figure_2.PNG}
                    \end{figure}
                    WLOG we can set $h=w$.

 Thus:

 
 \[lw+w^2+wl=27\]
 \[lw^2=23\]

 We can substitute $l$ and form a depressed cubic equation with $w$.\begin{align*}lw^2&=23\\l&=\frac{23}{w^2}\\2\left(\frac{23}{w^2}\right)w+w^2&=27\\\frac{46}{w}+w^2&=27\\w^2+\frac{46}{w}-27&=0\\w^3 -27w+46&=0.\end{align*}
Based on Rational Root Theorem the possible rational roots are $\pm1, \pm2, \pm23$

 A quick test reveals that $2$ is a root of the equation. Comparing coefficients we can factorize the equation into:

 $(w-2)(w^2+2w-23)=0$

 Besides $2$, we derive another positive root using the quadratic formula, $2\sqrt{6}-1$But to maximize the $m$ we need to pick the smaller $w$, which is $2$.

 Substituting this into $l=\frac{23}{w^2}$, we find that $l=\dfrac{23}4$.

 Applying it to our equation above:\begin{align*}4r^2&=(l+w+h)^2-54\\4r^2&=(l+2w)^2-54\\4r^2&=\left(\dfrac{23}4+2\cdot2\right)^2-54\\4r^2&=\left(\dfrac{39}4\right)^2-54\\4r^2&=\left(\dfrac{1521}{16}\right)-54\\4r^2&=\left(\dfrac{657}{16}\right)\\r^2&=\left(\dfrac{657}{64}\right). \end{align*}
$657+64=\boxed{721}$.

 





\textbf{Solution 3a (Derivative)}

to find the maximum m for $x^3-mx^2+27x-23=0$

 rewrite $m$ as function of $x$ and calculate derivatives to get maximum value, 
 \[m(x) =x + 27x^{-1} - 23x^{-2}\]
 \[m'(x) = 1 - 27x^{-2} -46x^{-3} = 0\]
 \[x^3 -27x+46=0\]
 \[(x-2)(x^2+2x-23)=0\]

 when $x = 2$, 
 \[m= 2 + \frac{27}{2}  - \frac{23}{4}  = \frac{39}{4}\]the rest is similar to solution 3

 





\textbf{Solution 4}

This question looks complex, but once converted into a number theory problem, it becomes elementary. We know, if the dimensions are taken to be numbers in the form of coprime numbers p/q,q/r, and r, it is immediately obvious that p=23.And solving we get: 
 \[23(r^2+q)/qr+q=27\] We know length cannot be -ve, in this case, therefore, q=4. And, again, we see:
 \[23/r+4+23r/4=27\] giving rise to r=2. For a cuboid inside a circle, we know:the radius is half its diagonal or, can we not say, 
 \[r^2=(L^2+b^2+h^2)/4\]or  here,
 \[r^2=(128+529)/64\]so here, 
 \[p+q=657+64=721\]

 





\textbf{Solution 5}

This problem essentially boils down to maximizing the value of $x+y+z$ (where $x$, $y$, and $z$ denote the dimensions of the box) given $xy+xz+yz=27$ and $xyz=23$. After doing so, we can calculate $r^2$ using $r^2=\frac{(x+y+z)^2-54}{4}$ (as in Solution 2). We can turn $x+y+z$ into an expression in terms of only $x$ and use the method of critical points. Since $xyz=23$, we have $yz=\frac{23}{x}$ and thus $xy+xz+yz=x(y+z)+\frac{23}{x}=27$. Isolating $y+z$, we find $y+z=\frac{27-\frac{23}{x}}{x}$, so
 \[x+y+z=x+\frac{27-\frac{23}{x}}{x}=x+\frac{27x-23}{x^2}=x+27\cdot x^{-1}-23\cdot x^{-2}.\]We have turned $x+y+z$ into an equation in $x$, using all conditions the problem has given us. We proceed with calculus. The maximum value of this function in $x$ is one of the critical points of this function, which can be calculated by equating the function's derivative with $0$. Using the power rule, the derivative of $f(x)=x+27\cdot x^{-1}-23\cdot x^{-2}$ is $f'(x)=1-27\cdot x^{-2}+46\cdot x^{-3}$. Equating $f'(x)$ to zero, we get
 \[1-27\cdot x^{-2}+46\cdot x^{-3}=0 \implies x^3-27x+46=0\]The roots of this can be found using the rational root theorem, yielding $x=2$ or $-1\pm 2\sqrt{6}$. Don't forget -- $x$ must be positive, so the only possible candidates of $x$ to maximize $f(x)$ are $2$ and $2\sqrt{6}-1$. Plugging both of them into $f(x)=x+27\cdot x^{-1}-23\cdot x^{-2}$, we find $x=2$ yields a greater result, $\frac{39}{4}$. Thus, the maximum value of $x+y+z=f(x)$ is $\frac{39}{4}$. Therefore, the value of $r^2$ is
 \[\frac{(\frac{39}{4})^2-54}{4}=\frac{657}{64} \implies p+q=\boxed{721}.\]

 





\textbf{Solution 6 (If you don't notice that two of the side lengths have to be equal)}

From the previous solutions we can see that we have to maximize $m$ in the expression $x^3 - mx^2 + 27x - 23 = 0$ such that $x^3 - mx^2 + 27x - 23 = 0$ has $3$ real and positive roots since the roots correspond to the side lengths of the rectangle. To do this we can use synthetic division to divide this expression by $(x - r)$ for some arbitrary root $r$ to try and see how $m$ relates to the roots. After doing synthetic division, we get 
 \[(x^2 + (r - m)x + (27 + r^2 - rm) + \frac{r^3 - r^2m + 27r - 23}{x - r})(x - r) = x^3 - mx^2 + 27x - 23\]Since we defined r to be a root of the expression, it must divide into it with no remainder, meaning that $r^3-r^2m+27r-23=0$ or $m = r + \frac{27}{r} - \frac{23}{r^2}$ 

 We can now substitute this value into $x^2 + (r-m)x + (27 + r^2 - rm)=0$ to get 

 $x^2 + (\frac{23}{r^2} - \frac{27}{r})x + \frac{23}{r} = 0$.

 Multiplying everything by $r^2$ yields $r^2x^2 + (23-27r)x + 23r = 0$. 

 We now have to find the maximum value of $r$ for which this expression has 2 roots, since $m$ strictly increases with $r$ from our previous expression of $m = r + \frac{27}{r} - \frac{23}{r^2}$

 To do this, we can find the discriminant of $r^2x^2 + (23-27r)x + 23r = 0$ and find the maximum value of $r$ for which it is positive. 

 The discriminant of this expression is $-92r^3 + 729r^2 - 1232r + 529$, and after using the rational root theorem we can find $\frac{23}{4}$ as a root. 

 The remaining expression is $(x - \frac{23}{4})(-92x^2 + 200x - 92)$. We can quickly see using the quadratic formula that $\frac{23}{4}$ is the largest root, meaning that $r=\frac{23}{4}$ is the largest value that a root of $x^3 - mx^2 + 27x - 23 = 0$ can have for which $x^3 - mx^2 + 27x - 23=0$ has $3$ positive, real solutions. From here, we can substitute this value back into the expression for $m$, giving us $m = \frac{39}{4}$. The length of the radius squared is $\frac{m^2-54}{4}=\frac{657}{64}$ and $657+64=\boxed{721}$ 

 




\end{enumerate}